% !TeX program = lualatex
% 共通の組版設定。本文と証明は各論文の節ファイルに一度だけ置く。
\documentclass[11pt,a4paper]{ltjsarticle}
\usepackage[margin=24mm]{geometry}
\usepackage{amsmath,amssymb,amsthm,mathtools,booktabs,longtable}
\usepackage{mathrsfs}
\usepackage[haranoaji]{luatexja-preset}
\usepackage[unicode,hidelinks]{hyperref}
\setlength{\emergencystretch}{3em}
\hypersetup{bookmarksnumbered=true}
\pdfstringdefDisableCommands{\def\qgr{qgr}\def\HH{HH}\def\Perf{Perf}}
\providecommand{\PublicationVersion}{publication-v32}
\newtheorem{theorem}{定理}[section]
\newtheorem{proposition}[theorem]{命題}
\newtheorem{lemma}[theorem]{補題}
\newtheorem{corollary}[theorem]{系}
\theoremstyle{definition}\newtheorem{definition}[theorem]{定義}
\newtheorem{assumption}[theorem]{入力・仮定}
\theoremstyle{remark}\newtheorem{remark}[theorem]{注意}
\renewcommand{\proofname}{証明}
\newcommand{\HH}{\operatorname{HH}}
\newcommand{\qgr}{\operatorname{qgr}}
\newcommand{\coh}{\operatorname{coh}}
\newcommand{\Perf}{\operatorname{Perf}}
\newcommand{\Ext}{\operatorname{Ext}}
\newcommand{\Hom}{\operatorname{Hom}}
\newcommand{\Spec}{\operatorname{Spec}}
\newcommand{\Spf}{\operatorname{Spf}}
\newcommand{\GL}{\operatorname{GL}}
\newcommand{\gldim}{\operatorname{gldim}}
\newcommand{\OO}{\mathcal O}
\renewcommand{\AA}{\mathbb A}
\newcommand{\PP}{\mathbb P}

% 各原稿の引用だけを共通文献データから組版する。
\makeatletter
\newcommand{\UsePublicationReferences}[1]{%
  \@for\pub@reference:=#1\do{%
    \expandafter\def\csname pubref@\pub@reference\endcsname{}%
  }%
}
\makeatother

\title{量子三次曲面のFricke幾何と周期\\
\large Poisson臨界被覆、$D_4$格子とtriality}
\author{}
\date{2026年11月18日\quad publication-v32}
\begin{document}
\maketitle
\begin{abstract}
三乗根量子 $\mathfrak{so}_3$ 包絡の中心と一次元表現から、Fricke三次曲面の幾何を調べる。
Askey--Wilson表示とCasimirの規約を明示し、量子化パラメータの微分が誘導する
Poisson括弧の係数を決定する。
相対Poisson臨界スキームと一次元表現スキームはそれぞれ次数15と5で有限平坦であり、
臨界値多項式は $N_5^2M_5$ と分解する。二つの五次被覆の同時幾何学的モノドロミーは
$S_5\times S_5$ である。
同じ表現被覆をLandau--Ginzburg臨界点族として記述し、
Cayley分岐曲面、局所Milnor格子とPicard格子の $D_4$、および四つの指数化周期を結ぶ。
対称族ではtrialityの不変・余不変格子の指数3が周期torusのisogenyを与え、
既知の $G_2$ 三成分写像と両立する。
具体的な座標比較と、引用する一般理論および未構成の圏同値を区別する。
\end{abstract}

\section{序論}\label{sec:fricke-introduction}
中心と表現スキームは、同じ非可換代数に付随する異なる可換幾何である。
本稿は三乗根量子 $\mathfrak{so}_3$ 包絡を具体例とし、
中心のFricke表示、表現被覆の分岐、コンパクト化の整数格子と周期を順に比較する。
以下の代数計算は標数0の代数閉体上で行い、周期と混合Hodge構造では基礎体を $\mathbb C$ とする。

\subsection{二つの五重被覆と四つの周期}
第一の中心結果は、単なる判別式の零集合でなく、その重複度と被覆を記述することである。
第\ref{sec:quantum-center-data}節の中心方程式を $P_{\mathbf a}(X,Y,Z,\Omega)=0$ とし、
\[
 \mathcal K=\Spec k[a,b,c,X,Y,Z,\Omega]/(P_{\mathbf a},P_X,P_Y,P_Z)
\]
を取る。一方、一次元表現は
\[
 \mathcal E=\Spec k[a,b,c,x,y,z]/
 ((1-q)yz-x-a,\ (1-q)zx-y-b,\ (1-q)xy-z-c)
\]
で与えられる。
\begin{theorem}[相対臨界被覆：第\ref{sec:relative-poisson-cover}節]
$\mathcal K$ と $\mathcal E$ は係数底 $\AA^3_{a,b,c}$ 上有限平坦で、
次数はそれぞれ15と5である。
$\Delta_{15}(D)=\det(D+\Omega\mid\mathcal O(\mathcal K))$、
$N_5(D)=\det(D+\Omega\mid\mathcal O(\mathcal E))$ とすると、
全係数で $\Delta_{15}=N_5^2M_5$、$M_5$ はmonicな五次多項式である。
二つの五重被覆が有限\'etaleである共通開集合上で、
その同時幾何学的モノドロミーは $S_5\times S_5$ となる。
\end{theorem}
有限平坦性は無限遠を含む完全交叉から、二重性は相対節点の局所模型から示す。
対称群の結論には、計算による分解型に加えて本文の群論と分岐の議論を用いる。

第二の中心結果は、Fricke型等位面
\[
 U_{A,B,C;h}=\{uvw-\tfrac12(u^2+v^2+w^2)-Au-Bv-Cw=h\}
\]
の四つの内部周期である。
滑らかな三次閉包 $S$ と境界三角形 $D=D_1+D_2+D_3$ を取ると、
$D^\perp\subset\operatorname{Pic}(S)$ は $D_4(-1)$ となる。
第\ref{sec:four-D4-periods}節のconic開集合と有限標識被覆上では、
四つの整数根基底について
\[
 \bigl(\mathcal P(\alpha_1),\ldots,\mathcal P(\alpha_4)\bigr)
 =\left(z_1/z_2,\ z_2/z_3,\ z_3/z_4,\ z_3z_4\right)
\]
を得る（定理\ref{thm:four-integer-D4-periods}）。
全 $z_i$ の捩れ性は、この範囲で有理混合Hodge拡大の分裂と同値である
（定理\ref{thm:four-period-rational-split}）。

対称族ではtrialityの不変格子から余不変格子への写像が指数3であり、
周期torusには $(\xi,\eta)\mapsto(\xi,\eta^3)$ が現れる。
そのWeyl商はBoalch--Palubaの既知の三成分写像と一致する
（定理\ref{thm:G2-period-isogeny-quotient}）。
共通の明示開集合で三成分の有理混合Hodge構造を比較できるが、
整数拡大類は一致せず、Kleinの例ではその位数が $7,21,21$ となる。

\subsection{先行研究との関係}
Fricke三次式と四穴球面のトレース座標はGoldman \cite{GoldmanTrace}、
量子的な生成関係とmodular作用はTerwilliger \cite{TerwilligerDAHA} に依存する。
本稿ではこれらの規約を三乗根包絡へ移し、Casimirのスカラー、
係数写像、およびPoisson括弧の係数9を直接計算する。
Cayley曲面のトーラス商とVieta力学はCantat--Loray \cite{CantatLorayCayley} の既知の理論である。
ここでは、それが五重の一次元表現被覆の分岐源になることを具体式で結ぶ。

三次曲面のホモロジー作用はGoldman--Neumann \cite{BoundaryGoldmanNeumann}、
線束制限と積分の周期比較はFriedman \cite{FrickeFriedman}、
対称Fricke族の $G_2$ 指標多様体としての実現はBoalch--Paluba \cite{FrickeBoalchPaluba}
を用いる。これらの一般理論を新規結果として主張しない。
本稿の役割は、この一つの表現被覆における整数格子・係数・周期の比較を
写像と式で与えることであり、個々の式や帰結の網羅的な既出性比較はなお必要である。
特に、次数3の底写像を群そのものの中心isogenyや非可換圏の同値と読み替えない。

\subsection{適用範囲}
大域的な非可換三次曲面の形式変形とH3認識は関連論文\cite{CompanionE6}で扱う。
本稿の中心・臨界被覆・周期の主結果は、その全モジュライ分類を仮定しない。
H3への言及は、構造対象を保つexact同値と飽和反標準環の認識が成り立つ場合に限定する。
spherical DAHAとのMorita同値、非可換H3圏上のVieta自己同値、
Fukaya圏との同値は本稿の結論に含めない。
周期に見えないVieta力学の境界での記述は\cite{CompanionBoundary}に委ねる。

\tableofcontents
\section{基本代数と二種類の座標}\label{sec:quantum-center-data}
$q^2+q+1=0$ とし、$\mathbf a=(a,b,c)$ に対して
\[
 V_{\mathbf a}=k\langle x,y,z\rangle/
 (yz-qzy-x-a,\ zx-qxz-y-b,\ xy-qyx-z-c)
\]
と置く。そのRees環を $B_{\mathbf a}$、Rees変数を $t$ と書く。
PBW基底と全中心式の証明は\cite[三乗根包絡と中心三次切断]{CompanionE6}にある。
ここでは後続の計算に必要な全係数を明記する：
\begin{align}
 X&=x^3+q^2x,\quad Y=y^3+q^2y,\quad Z=z^3+q^2z,\notag\\
 \Omega&=xyz+\frac{q-1}{3}(x^2+z^2)+\frac{q+2}{3}y^2
 +(q-1)^{-1}(ax+qby+cz).\label{eq:Omega}
\end{align}
斉次化した各三次元には添字 $h$ を付ける。
$f=(1-q)/3$、$e=-(2q+1)/9$ とし、
\begin{align*}
 \ell_a&=((2q+1)a^3+2(q-1)a+3qbc)/9,\quad \ell_b,\ell_c\text{ は巡回置換},\\
 K&=((q+1)(a^2+b^2+c^2)+(q+2)abc)/3,\\
 L&=((1-q)(a^2b^2+a^2c^2+b^2c^2)+(q+2)(a^2+b^2+c^2))/9
\end{align*}
と置けば、
\begin{equation}\label{eq:center}
\begin{split}
 P_{\mathbf a}={}&\Omega^3+f\Omega^2-XYZ+e(X^2+Y^2+Z^2)\\
 &-(aX+bY+cZ)\Omega+\ell_aX+\ell_bY+\ell_cZ+K\Omega+L=0
\end{split}
\end{equation}
が全中心の一本の関係である。
この全係数式は計算検証器でも原関係から再現できる。
Casimir切断は $\Omega=-D$ とする。

前半のFricke中心座標と、後半の一次元表現のLG座標には、
どちらにも $u,v,w$ が現れるが意味は異なる。
第\ref{sec:daha-character}節では $(-X/e,-Y/e,-Z/e)$、
第\ref{sec:cayley-lg-connections}節以後では $(1-q)(x,y,z)$ を用いる。
後者の係数は $(A,B,C)=(1-q)(a,b,c)$ であり、
LG level $h$ と中心切断値 $D$ も別の量である。
各比較写像でこれらを明示する。

\section{DAHA・四穴球面character variety・Goldman幾何への厳密な接続}
\label{sec:daha-character}
本節では三捩れの量子 $\mathfrak{so}_3$ 包絡の\emph{特別な}
Casimir切断を扱う。一般の中央三次切断全体や、楕円境界を動かす八次元の
量子化族を同一のcharacter varietyと同定する主張ではない。
基礎体は標数零で原始三乗根 $q$ を含むとし、
$V_{a,b,c}=B_{\mathbf a}/(t-1)$、$\Omega=\Omega_{\mathbf a,h}|_{t=1}$ とする。

\begin{theorem}[普遍Askey--Wilson代数とCasimirの完全な対応]
\label{thm:aw-casimir-exact}
Terwilligerの普遍Askey--Wilson代数 $\Delta_q$ の生成元を
$A_{\rm AW},B_{\rm AW},C_{\rm AW}$、その三つの中心パラメータを
$\alpha_{\rm AW},\beta_{\rm AW},\gamma_{\rm AW}$ と書く。
\cite[Definitions 2.1, 2.2, 2.6]{TerwilligerDAHA} の規約を使うと
\[
 V_{a,b,c}\cong\Delta_q/
 (\alpha_{\rm AW}-(1-q)a,\,
 \beta_{\rm AW}-(1-q)b,\,
 \gamma_{\rm AW}-(1-q)c),
 \qquad (A_{\rm AW},B_{\rm AW},C_{\rm AW})=(1-q)(x,y,z).
\]
この同型の下でCasimir元は\emph{厳密に}
\begin{equation}\label{eq:aw-casimir-scale}
 \Omega_{\rm AW}=q(1-q)^3\Omega=3(q+2)\Omega
\end{equation}
となる。従って任意の $D\in k$ について
\[
 V_{a,b,c}/(\Omega+D)\cong
 \Delta_q/(\alpha_{\rm AW}-(1-q)a,\,
 \beta_{\rm AW}-(1-q)b,\,
 \gamma_{\rm AW}-(1-q)c,\,
 \Omega_{\rm AW}+3(q+2)D).
\]
これらは代数の表示同型であり、spherical DAHAとのMorita同値は含まない。
\end{theorem}
\begin{proof}
Terwilligerの三つの中心関係は $q^3=1$ のとき
$B_{\rm AW}C_{\rm AW}-qC_{\rm AW}B_{\rm AW}
 =(1-q)(A_{\rm AW}+\alpha_{\rm AW})$ とその巡回式になる。
表示を比較すれば最初の同型を得る。
Casimirの原式は
\[
 \Omega_{\rm AW}=q^{-1}A_{\rm AW}C_{\rm AW}B_{\rm AW}
 +q^{-2}A_{\rm AW}^2+q^{-2}B_{\rm AW}^2+q^2C_{\rm AW}^2
 -q^{-1}A_{\rm AW}\alpha_{\rm AW}
 -q^{-1}B_{\rm AW}\beta_{\rm AW}-qC_{\rm AW}\gamma_{\rm AW}.
\]
これに各元のスケールを代入し、$zy=q^2yz-q^2(x+a)$ によって
$xzy$ をPBW順に直すと
\[
 \Omega_{\rm AW}=(1-q)^2
 \bigl(q(1-q)xyz+q^2(x^2+z^2)+qy^2
      -qax-q^2by-qcz\bigr).
\]
$q^2+q+1=0$ を使うと右辺は
$q(1-q)^3[xyz+(q-1)(x^2+z^2)/3+(q+2)y^2/3
 -(q+2)(ax+qby+cz)/3]$ に一致する。
これは式\eqref{eq:Omega}の非斉次化である。
最後に $q(1-q)^3=3(q+2)$ が従う。
式は \texttt{verify\_daha\_character.py} で独立に厳密検算する。
\end{proof}

\subsection{Fricke三次曲面と四穴球面のトレース}
\begin{theorem}[character variety係数空間への優越写像]
\label{thm:fricke-trace-dominant}
$f=(1-q)/3$、$e=-(2q+1)/9$ とし、式\eqref{eq:center}の
$\ell_a,\ell_b,\ell_c,K,L$ を用いる。
$\Omega=-D$ とおき、中心の残りの生成元を
$u=-X/e,\ v=-Y/e,\ w=-Z/e$ と置くと、中心は
\begin{equation}\label{eq:fricke-e6}
 F=u^2+v^2+w^2+uvw+A u+B v+C w+T=0
\end{equation}
で与えられる。ただし
\begin{align}\label{eq:fricke-e6-parameters}
 A&=-\frac{Da+\ell_a}{e^2},&
 B&=-\frac{Db+\ell_b}{e^2},&
 C&=-\frac{Dc+\ell_c}{e^2},\\
 T&=-\frac{D^3-fD^2+DK-L}{e^3}.\notag
\end{align}
係数写像 $\Psi:\AA^4_{a,b,c,D}\to\AA^4_{A,B,C,T}$ はdominantかつ
generically finiteであり、$a=b=c=0$ におけるJacobianは
\begin{equation}\label{eq:fricke-jacobian}
 \det(d\Psi)=\frac{3}{e^9}\,D\bigl(D-2f/3\bigr)^4.
\end{equation}

四穴球面 $\Sigma_{0,4}$ の境界monodromyのトレースを
$r_1,r_2,r_3,r_4$ とすれば、その相対 $\mathrm{SL}_2$ character varietyの
Fricke三次式は\eqref{eq:fricke-e6}で、係数は
\begin{equation}\label{eq:fricke-trace-map}
\begin{split}
 A&=-(r_1r_2+r_3r_4),\qquad
 B=-(r_2r_3+r_1r_4),\\
 C&=-(r_1r_3+r_2r_4),\qquad
 T=\sum_{i=1}^{4}r_i^2+r_1r_2r_3r_4-4.
\end{split}
\end{equation}
このトレース係数写像もdominantかつgenerically finiteである。
従って適切な稠密開集合上で有限分離拡大（または有限\'{e}tale底変更）を
行えば、$V_{a,b,c}/(\Omega+D)$ の中心を
$\Sigma_{0,4}$ の相対character varietyとして同定できる。
これは中心の同型であり、元の非可換代数が
character stackそのものになるという主張ではない。
\end{theorem}
\begin{proof}
中心関係\eqref{eq:center}に $\Omega=-D$ を代入して符号を反転すると
\[
 XYZ-e(X^2+Y^2+Z^2)
 -(Da+\ell_a)X-(Db+\ell_b)Y-(Dc+\ell_c)Z
 +D^3-fD^2+DK-L=0.
\]
$X=-eu,Y=-ev,Z=-ew$ を代入して $-e^3$ で割れば
\eqref{eq:fricke-e6}--\eqref{eq:fricke-e6-parameters} が出る。
$\partial\ell_a/\partial a|_{\mathbf a=0}=2(q-1)/9=-2f/3$ で、
$K,L$ は $(a,b,c)$ に一次項を持たない。
従ってこの点でJacobianは対角で、対角項は
$-(D-2f/3)/e^2$ が三回と $-3D(D-2f/3)/e^3$ が一回である。
式\eqref{eq:fricke-jacobian} が従い、右辺は零多項式でない。
等式\eqref{eq:fricke-trace-map} は
Goldman \cite[\S5.2.1, equation (5.2.1.2)]{GoldmanTrace} の四穴球面
トレース関係である。
その四変数Jacobian行列式を $(r_1,r_2,r_3,r_4)=(1,2,3,4)$ で評価すると
$-200$ となり、こちらもdominantかつgenerically finiteである。
両写像の一般の整次元零次fiberは有限分離であるから、
有限\'{e}tale開部分上のbase changeで同型の記述を得る。
\end{proof}

\subsection{Braid作用はFricke--Vieta変換を持ち上げる}
\begin{theorem}[三乗根での量子的Vieta変換]
\label{thm:aw-braid-vieta}
$a,b,c$ を可変の中央パラメータとみなした $V_{a,b,c}$ の全族には
Terwilligerの $\mathrm{PSL}_2(\mathbb Z)$ 作用
\cite[Lemmas 2.3, 2.5, 2.8]{TerwilligerDAHA} がある。
その生成元の一つ $\sigma$ は
\[
 \sigma(a,b,c)=(b,a,c),\qquad
 \sigma(x)=y,\quad\sigma(y)=x,\quad
 \sigma(z)=(q^2-q)xy-q^2z+qc
\]
で与えられ、Casimir $\Omega$ と中心切断値 $D$ を固定する。
中心の作用は
\begin{equation}\label{eq:vieta-center}
 \sigma(X)=Y,\quad \sigma(Y)=X,\quad
 \sigma(Z)=\frac{XY+c\Omega-\ell_c}{e}-Z.
\end{equation}
特にFricke座標と係数の上では
\begin{equation}\label{eq:vieta-fricke}
 (u,v,w;A,B,C,T)\longmapsto
 (v,u,-uv-C-w;B,A,C,T)
\end{equation}
となる。もう一つの生成元 $\rho$ は両方の三つ組を巡回置換し、$T$ を固定する。
これは中心上の通常のFricke--Vieta型変換の\emph{非可換持上げ}である。
一般の楕円境界型への同じ作用の延長は示していない。
\end{theorem}
\begin{proof}
Terwilligerの作用は $A_{\rm AW}\leftrightarrow B_{\rm AW}$、
$C_{\rm AW}\mapsto C_{\rm AW}+(A_{\rm AW}B_{\rm AW}-B_{\rm AW}A_{\rm AW})/(q-q^{-1})$
であり、三中心パラメータの前二つを交換しCasimirを固定する。
この式を $x,y,z$ へ移し、
$yx=q^2xy-q^2(z+c)$ を用いると、上の $\sigma(z)$ が得られる。
元の三関係が保たれることもPBW還元で直接確認できる。

$P$ を式\eqref{eq:center}の中心三次方程式とすると、
$a,b$ と $X,Y$ の同時交換で $P$ は不変であり、
$Z$ に関して
\[
 P=eZ^2+(-XY-c\Omega+\ell_c)Z+P_0(X,Y,\Omega)
\]
の形である。$\sigma(Z)$ も同じ二次式の根だから、整域性により
$\sigma(Z)=Z$ または式\eqref{eq:vieta-center}の第二根となる。
ところがPBW次数で $\sigma(z)$ の最高記号は $(q^2-q)xy$、
$\operatorname{gr}\sigma(Z)$ の最高記号は
$(q^2-q)^3x^3y^3=e^{-1}x^3y^3$ で次数6を持つ。
よって次数3の $Z$ とは異なり、後者の根が選ばれる。
さらに \texttt{verify\_daha\_braid.py} では、三つの変換後関係と
Casimir不変性および式\eqref{eq:vieta-center}を、
元のPBW規則で全係数検証する。
$u=-X/e$ 等を代入すれば式\eqref{eq:vieta-fricke}となる。
\end{proof}

\begin{proposition}[Goldman型Poisson構造と対数Calabi--Yauコンパクト化]
\label{prop:fricke-poisson-logcy}
各係数 $(A,B,C,T)$ に対する\eqref{eq:fricke-e6}の中心には
Jacobian Poisson括弧
\[
 \{u,v\}=\frac{\partial F}{\partial w}=2w+uv+C,\quad
 \{v,w\}=2u+vw+A,\quad
 \{w,u\}=2v+uw+B
\]
がある。これはGoldman括弧と符号・非零の全体定数の規約を除いて一致する
\cite{GoldmanTrace,HikamiCluster}。
式\eqref{eq:vieta-fricke}および巡回作用は係数も同時に変換する全族のPoisson同型である。

さらに射影閉包
\[
 \bar C=\{t(u^2+v^2+w^2)+uvw+t^2(Au+Bv+Cw)+Tt^3=0\}\subset\PP^3
\]
は無限遠 $t=0$ に三角形 $D_\infty=\{uvw=0\}$ を持ち、
$\bar C$ はこの境界に沿って常に滑らかである。
従ってaffine中心が滑らかなら $\bar C$ は滑らかな三次曲面で、
$D_\infty\in|-K_{\bar C}|$ は三本の直線からなる反標準cycleである。
$(\bar C,D_\infty)$ はLooijenga pairとなり、
\[
 \omega=\operatorname{Res}_{F=0}\frac{du\wedge dv\wedge dw}{F}
\]
はそのlog symplectic形を与える。
この形とPoisson括弧の比較は古典的な三次曲面の内容であり、
本研究の量子代数でのBrown--Gordon Poisson orderが
同じ正規化を持つことまでは本命題では証明しない。
\end{proposition}
\begin{proof}
$\nabla F$ を使って定める三変数のJacobian括弧はJacobi恒等式を満たし
$F$ をCasimirとするので、曲面に降りる。
Vieta置換 $w\mapsto-uv-C-w$ は $F$ を保存しJacobian determinantが $-1$ だから
反Poissonである。$u,v$ の交換も反Poissonなので、両者の合成
\eqref{eq:vieta-fricke}はPoissonとなる。
巡回置換はdeterminant $+1$ である。
Goldman括弧との照合は四穴球面のトレース座標の既知公式を用いる。

$t=0$ で射影方程式は $uvw=0$ となる。
その辺上で $u,v$ など二座標が非零なら残りの座標による偏微分が非零であり、
三頂点のそれぞれでは $t$ による偏微分が一つの座標の平方となり非零である。
ゆえに無限遠全体で滑らかである。
$K_{\bar C}=\OO_{\bar C}(-1)$ はadjunctionによるから、
超平面切断 $D_\infty$ は反標準であり、上のresidueが
$K_{\bar C}(D_\infty)\simeq\OO_{\bar C}$ を自明化する。
\end{proof}

\begin{proposition}[三角境界からの $D_4\subset E_6$ 根格子（既知の幾何）]
\label{prop:d4-inside-e6}
$\bar C$ が滑らかな三次曲面で $D_\infty=D_1+D_2+D_3$ が
三本の直線の三角形とする。このとき交叉形式について
\[
 K_{\bar C}^{\perp}\cong E_6(-1),\qquad
 \langle D_1,D_2,D_3\rangle^{\perp}\cong D_4(-1).
\]
特に、四穴球面のFricke三次曲面で現れる $D_4$ 根格子は、
三次del Pezzoの $E_6$ 根格子の、反標準三角形を固定した部分として具体化する。
これは古典的な三次曲面・Naruki--Looijenga型周期理論の既知結果
\cite{SchockD4} を本稿の座標に適用したものである。
\end{proposition}
\begin{proof}
三次曲面上のtritangent planeの標準的なmarkingを選べば、
$\bar C\simeq\operatorname{Bl}_{p_1,\ldots,p_6}\PP^2$ で、
直線classを $H$、例外曲線classを $E_i$ として
\[
 D_1=H-E_1-E_2,\quad D_2=H-E_3-E_4,\quad D_3=H-E_5-E_6.
\]
従って $D_1+D_2+D_3=3H-\sum E_i=-K_{\bar C}$。
$\operatorname{Pic}(\bar C)$ は
$H^2=1,E_i^2=-1,H\cdot E_i=0$ の交叉形式を持つ。
$L=aH-\sum b_iE_i$ が三本の $D_i$ と直交する条件は
$b_1+b_2=b_3+b_4=b_5+b_6=a$ である。
この整数解格子の基底は
\[
 \alpha_0=H-E_1-E_3-E_5,\quad
 \alpha_1=E_1-E_2,\quad
 \alpha_2=E_3-E_4,\quad
 \alpha_3=E_5-E_6
\]
であり、$\alpha_i^2=-2$、$\alpha_0\cdot\alpha_j=1$ ($j=1,2,3$)、
三つの外側rootは互いに直交する。
そのGram行列は負の$D_4$型Cartan行列なので結論を得る。
また $K^\perp\cong E_6(-1)$ は六点blow-upの標準根格子である。
整数基底とGram行列の検算は
\texttt{computations/verify\_d4\_lattice.py} にも記録する。
\end{proof}

\begin{remark}[先行研究との区別と次の比較]
Goldman--Toledo \cite{GoldmanToledo} はさらに、無限遠が非退化なtritangent planeである\emph{任意の}affine三次曲面が
四穴球面の相対character varietyとなることを示す。
従ってFricke型方程式への同定そのものは既知であり、
本研究での位置付けは指定した量子包絡の中心・Casimir規約を通じた
明示的な比較にある。
四穴球面skein代数と普遍Askey--Wilson代数の同定、rank-one DAHAへの写像、
$\mathcal U_q(\mathfrak{sl}_2)$ のRacah問題、Braid作用は
\cite{TerwilligerDAHA,AWAvatars,CookeLacabanne} で研究されている。
また四穴球面のcluster座標・Painlev\'{e} confluenceは \cite{HikamiCluster,CMR}、
三角境界のlog Calabi--Yau三次曲面のmirrorとskein代数は \cite{BousseauSkein} に現れる。
本稿の新しい計算は、\emph{指定した}三乗根・PBW規約の下での
Casimirスケール、Fricke係数写像、量子的Vietaの正規形を
明示し同一のH3族と接続することである。
これらの個別式の文献上の既出性は未確認であり、新規発見と断定しない。

特に、$\Omega=-D$ は $\beta=\gamma=\delta=0$ の三角境界切断であり、
一般の滑らかな楕円境界の七次元正常形全体ではない。
$\widetilde D_4$ 型の四穴球面DAHAと、楕円 $\widetilde E_6$ 特異点からの
Etingof--Ginzburg量子化 \cite{EG}、GDAHAの $\widetilde E_6$ 枝 \cite{EOR}
を同一の代数的構成と混同してはならない。
今後は有限位数でのskein--DAHA同型の基底変更の正確な適用範囲、
Frobenius中心とPoisson order、Vieta--Braid作用の相対H3心への延長、
七・八次元の楕円境界への変形とcluster/mirror座標を比較する。
\end{remark}
\section{三乗根Poisson構造と、楕円切断・H3心への二つの障害}
\label{sec:root-poisson-orders}
本節は第\ref{sec:daha-character}節の中心の幾何学を、
非可換代数そのものの変形と結ぶ。日付は2026年11月17日、内容版はv15である。
三乗根における量子化からGoldman構造が現れる一般的機構は既知である。
Baseilhac--Rocheの一次論文\footnote{S. Baseilhac, P. Roche,
\emph{Unrestricted Quantum Moduli Algebras. I. The Case of Punctured Spheres},
SIGMA 18 (2022), 025, 78 pp.,
\href{https://doi.org/10.3842/SIGMA.2022.025}{doi:10.3842/SIGMA.2022.025},
Theorems 7.17, 7.19 and Corollary 7.20(3).}
はFrobenius写像を通じたこの関係を証明する。
本節で独立に証明するのは、本稿のPBW規約での正確な係数、
中心三次切断をその変形に持ち上げる際の障害、
および量子的Vieta作用と固定H3心の非両立性である。
個別の表示式や障害の文献上の新規性は未確定である。

\subsection{量子化パラメータの変形と正確な係数}
$q^2+q+1=0$ とし、$a,b,c$ は固定する。
$Q=q\exp(\hbar)$ とおき、$\hbar$-進完備化した代数
\[
 \mathcal V_\hbar=k[[\hbar]]\langle x,y,z\rangle/
 (yz-Qzy-x-a,\ zx-Qxz-y-b,\ xy-Qyx-z-c)
\]
を考える。$\hbar$ と斉次化変数 $t$ は異なる。
ここでの変形は、一般のSklyanin方向を同時に動かす変形ではない。

\begin{lemma}[全係数PBWと持ち上がるCasimir]
\label{lem:generic-q-casimir}
$Q$ を不定元としても上の三関係はPBW基底 $x^iy^jz^k$ を与え、
\begin{equation}\label{eq:generic-q-casimir}
 \Omega_Q=xyz+\frac{x^2+Q^2y^2+z^2}{Q^2-1}
       +\frac{ax+Qby+cz}{Q-1}
\end{equation}
は中心元である。$Q=q$ で $\Omega_Q$ は本稿の $\Omega$ に一致する。
任意の $D$ に対し、$\mathcal V_\hbar/(\Omega_Q+D)$ も
$k[[\hbar]]$ 上平坦である。
\end{lemma}
\begin{proof}
書換え規則は
\[
 yx=Q^{-1}xy-Q^{-1}(z+c),\quad
 zx=Qxz+y+b,\quad
 zy=Q^{-1}yz-Q^{-1}(x+a).
\]
全次数、次いで逆順対を下げる順序で停止し、唯一の重なり $zyx$ の両還元は一致する。
実際、両方とも
\[
 Q^{-1}xyz+Q^{-1}(y^2+by-x^2-ax-z^2-cz)
\]
となる。式\eqref{eq:generic-q-casimir}の三交換子に同じ規則を使うと全て零となる。
これは $\mathbb Q(Q)[a,b,c]$ 上の恒等式であり、有限標本による外挿ではない。
$q^2-1\ne0$ なので形式底への代入は正則で、特殊値は式\eqref{eq:Omega}となる。
PBWフィルトレーションの先頭環は三変数量子多項式環であり、
$\operatorname{gr}(\Omega_Q+D)=xyz$ は正則中心元である。
従って商の随伴次数環はこの量子多項式環の $(xyz)$ 商である。
少なくとも一つの指数が零のPBW語が商の基底となり、平坦性を得る。
形式完備化ではこの基底に対応する位相的自由加群を用いる。
\end{proof}

\begin{theorem}[root-of-unity括弧はGoldman型括弧の9倍]
\label{thm:root-poisson-goldman}
$V=V_{a,b,c}$ とし、その全中心を
$Z(V)=k[X,Y,Z,\Omega]/(P)$ と書く。$P$ は式\eqref{eq:center}の左辺である。
中心元 $f,g$ の持上げ $\widetilde f,\widetilde g$ に対し
\[
 \{f,g\}_{\rm rt}=\left.\frac{[\widetilde f,\widetilde g]}{\hbar}\right|_{\hbar=0}
\]
は持上げによらないPoisson括弧を定める。その値は
\begin{equation}\label{eq:root-poisson-center}
 \begin{aligned}
 \{X,Y\}_{\rm rt}&=-9P_Z=9(XY-2eZ+c\Omega-\ell_c),\\
 \{Y,Z\}_{\rm rt}&=-9P_X,\qquad
 \{Z,X\}_{\rm rt}=-9P_Y,\qquad
 \{\Omega,-\}_{\rm rt}=0.
 \end{aligned}
\end{equation}
$a,b,c$ を変数として扱う場合、これらもPoisson中心に属する。
特に $\Omega=-D$ の切断上のFricke座標では
\begin{equation}\label{eq:root-goldman-nine}
 \{u,v\}_{\rm rt}=9(2w+uv+C),\quad
 \{v,w\}_{\rm rt}=9(2u+vw+A),\quad
 \{w,u\}_{\rm rt}=9(2v+uw+B).
\end{equation}
従って第\ref{sec:daha-character}節で固定したJacobian規約の
Goldman型括弧の正確に9倍である。
\end{theorem}
\begin{proof}
PBW平坦性から交換子を $\hbar$ で割れる。
二つの持上げの差 $\hbar f_1$ を入れたときの差は
$[f_1,g]$ の特殊値であり、$g$ が中心なので零である。
Jacobi恒等式により得られる元は中心に入り、Leibniz則とJacobi則も交換子から従う。
$\Omega_Q$ とパラメータが厳密に中心へ持ち上がることから最後の等式を得る。

正確な係数は次の全係数計算で決まる。
$\widetilde X=x^3+Q^2x$, $\widetilde Y=y^3+Q^2y$, $\widetilde Z=z^3+Q^2z$ と置き、
上の三規則で積をPBW正規形 $\operatorname{NF}_Q$ に直すと
\begin{align*}
 \left.(Q\partial_Q)\operatorname{NF}_Q
   [x^3+Q^2x,y^3+Q^2y]\right|_{Q=q}
  ={}&9\bigl(x^3y^3+q^2x^3y+q^2xy^3+qxy\bigr)\\
   &-18e(z^3+q^2z)+9c\Omega-9\ell_c.
\end{align*}
右辺の最初の括弧は $XY$ であり、求める等式となる。
残る二つは巡回置換で得る。この等式は
\texttt{computations/verify\_root\_poisson.py} が
$\mathbb Q(Q)[a,b,c]$ 上で積を形成し、微分後に
$\mathbb Q(\sqrt{-3})[a,b,c]$ へ厳密代入して、全PBW係数を零と確認する。
また同じコードは三つのCasimir交換子と九つの特殊中心交換子も別に確認する。
$Q\partial_Q$ は選んだ変形 $Q=q\exp(\hbar)$ の一次微分に等しい。

最後に $P|_{\Omega=-D,X=-eu,Y=-ev,Z=-ew}=e^3F$ なので
$P_Z=-e^2F_w$、$\{u,v\}=e^{-2}\{X,Y\}=9F_w$ である。
これは括弧を後から便宜上定義したものではなく、元代数の変形からの同定である。
\end{proof}

\begin{remark}[規約と先行研究]
変形を $Q=q+\hbar$ と取れば係数は $9q^{-1}$ になる。
従って数値9を変形パラメータの選択から切り離してはならない。
Huang\footnote{H.-W. Huang, \emph{Center of the universal Askey--Wilson algebra
at roots of unity}, Nuclear Phys. B 909 (2016), 260--296,
\href{https://doi.org/10.1016/j.nuclphysb.2016.05.006}{doi:10.1016/j.nuclphysb.2016.05.006};
\href{https://arxiv.org/abs/1509.08719}{arXiv:1509.08719}.}
によるroot-of-unity中心の一般表示と、本稿の三乗根の明示計算は区別する。
skein/DAHA同型が全ての中央特殊化に保たれることを、本定理の前提にはしていない。
\end{remark}

\subsection{symplectic leaf上の表現論}
\begin{corollary}[Fricke切断のPoisson orderと中心fiber]
\label{cor:root-poisson-fibers}
$k=\mathbb C$ とする。$V$ は上の中心上のPoisson orderである。
$A_D=V/(\Omega+D)$ とし、$C_D=\Spec(k[X,Y,Z,\Omega]/(P,\Omega+D))$ と置く。
この中心部分環は $A_D$ に単射し、$A_D$ はその上有限なPoisson orderとなる。
$C_D$ のsmooth locusは一つのsymplectic leafであり、各特異点は零次元leafである。
従って $p,p'\in C_D^{\rm sm}$ について
\[
 A_D/\mathfrak m_pA_D\cong A_D/\mathfrak m_{p'}A_D
\]
が成り立つ。一点でこの代数が $M_3(\mathbb C)$ なら、smooth locusの全点で同じである。
この最後の条件を、全 $D$ について追加検証なしに主張しない。
\end{corollary}
\begin{proof}
$x^3=X-q^2x$ と巡回式により、$V$ は $k[X,Y,Z]$ 上
$0\le i,j,k<3$ の27語で自由である。独立性は最高PBW記号で分かる。
従って全中心上有限であり、そのgeneric rankは $27/3=9$ である。
中心元の線型な持上げを選び、同じ交換子微分を $V$ 全体へ作用させると導分を得る。
選択を変えた際の差は内部導分で、中心上の括弧は変わらない。
これはBrown--GordonのPoisson orderの定義と量子化構成に一致する。\footnote{
K. A. Brown, I. Gordon, \emph{Poisson orders, symplectic reflection algebras and
representation theory}, J. Reine Angew. Math. 559 (2003), 193--216,
\href{https://arxiv.org/abs/math/0201042}{arXiv:math/0201042},
\S\S2.1--2.2 and Theorem 4.2. ここでは同原論文の定義を使い、
中心元を全代数の導分へ送る写像の選択非依存性は主張しない。}
$\Omega_Q+D$ が厳密に中心なので導分は商へ降りる。
$V$ は整域であり、$z=(\Omega+D)v$ が中心なら
$(\Omega+D)[v,a]=0$ から $v$ も中心である。
よって $Z(V)\cap(\Omega+D)V=(\Omega+D)Z(V)$ であり、中心部分環の単射を得る。

Fricke方程式は全パラメータで既約である。
もし因数分解すれば、全次数の先頭式 $uvw$ により一次因子の先頭は
$u,v,w$ のいずれかの倍数である。
例えば $u-u_0$ を因子に持つとすれば、$u=u_0$ を代入した式の $v^2$ 係数1に矛盾する。
射影閉包は無限遠の三角形に沿って滑らかである。
特異曲線があれば射影閉包内で無限遠に交わるため矛盾し、特異点は有限である。
従ってsmooth locusは連結であり、そこで式\eqref{eq:root-goldman-nine}は非退化、
特異点では零となる。各leafは局所閉で、これらはsymplectic coreにも一致する。
Brown--Gordon Theorem 4.2を適用して中心fiberの同型を得る。
$M_3$ というサイズはgeneric rank 9と整合するが、
特定のCasimir値がgeneric Azumaya開集合と交わるかは別の確認事項である。
\end{proof}

\subsection{純粋な量子化パラメータ変形が通れない中心切断}
\begin{theorem}[非Casimir切断の一次平坦持上げ障害]
\label{thm:poisson-cut-obstruction}
$G=\Omega+\beta X+\gamma Y+\delta Z+\varepsilon\in Z(V)$ とする。
次の条件は同値である。
\begin{enumerate}
\item $(G)\subset Z(V)$ が $\{-,-\}_{\rm rt}$ に関するPoissonイデアルである。
\item $\beta=\gamma=\delta=0$。
\item $V/(G)$ が、固定した変形 $\mathcal V_\hbar$ の
$k[[\hbar]]$-平坦な両側イデアル商へ持ち上がる。
\end{enumerate}
(3)を $k[\hbar]/(\hbar^2)$ 上の平坦な両側イデアル商と弱めても同値である。
これは指定したambient変形の内部での障害であり、抽象的な圏変形や
Sklyanin方向を同時に動かす変形の非存在を意味しない。
\end{theorem}
\begin{proof}
$G=0$ で $\Omega$ を消去する。
残った中心関係 $P_G\in k[X,Y,Z]$ はちょうど次数3である。
その最高斉次部は
$-(\beta X+\gamma Y+\delta Z)^3-XYZ$ であり、恒等的に零とはならない。
従って次数2以下の多項式が $k[X,Y,Z]/(P_G)$ で零になることと、
その多項式自身が零であることは同値である。
式\eqref{eq:root-poisson-center}より、$G=0$ で計算した
$\{G,X\}/9,\{G,Y\}/9,\{G,Z\}/9$ の二次斉次部はそれぞれ
\begin{equation}\label{eq:cut-leading-obstructions}
 X(\delta Z-\gamma Y),\qquad
 Y(\beta X-\delta Z),\qquad
 Z(\gamma Y-\beta X).
\end{equation}
これらが全て零なら $\beta=\gamma=\delta=0$ である。
逆は $\Omega$ のCasimir性から明らかで、(1)と(2)を得る。

必要な一般補題を直接示す。
平坦変形 $\mathcal A_\hbar$ の平坦商 $\mathcal A_\hbar/\mathcal I_\hbar$ があり、
特殊イデアルを $I$ とする。$f\in I\cap Z(\mathcal A_0)$ と
$z\in Z(\mathcal A_0)$ の持上げを取ると
$[\widetilde f,\widetilde z]\in\mathcal I_\hbar\cap\hbar\mathcal A_\hbar$ である。
商の平坦性からこの交叉は $\hbar\mathcal I_\hbar$ に等しい。
従って $\{f,z\}_{\rm rt}\in I\cap Z(\mathcal A_0)$ となる。
同じ議論は双対数底にも有効である。
$V$ の整域性から $(G)V\cap Z(V)=(G)Z(V)$ なので、(3)は(1)を必要とする。
(2)なら補題\ref{lem:generic-q-casimir}の商が全次数の持上げを与える。
\end{proof}

\begin{corollary}[本稿の選択点における具体的障害]
$a=b=c=0$、$G_0=2\Omega+X+2Y+3Z+5$ の点では、
$\{G_0,X\}_{\rm rt}\bmod G_0$ の二次斉次部は
$9X(3Z-2Y)\ne0$ である。
従ってこの切断は純粋な $Q=q\exp(\hbar)$ 変形の平坦商としては
一次にも持ち上がらない。$a,b,c$ の一次補正を同時に許しても、この障害は変わらない。
\end{corollary}
\begin{proof}
前定理を $G_0/2$ に適用する。
さらに中心元同士の交換子は $Q=q$ において全 $a,b,c$ で零なので、
$a,b,c$ によるその一次微分も零である。
従って低次係数の変形を加えても中心括弧の障害を消せない。
\end{proof}

\subsection{affine modular対称性と反標準コンパクト化の相違}
\begin{proposition}[量子的Vietaは固定pointed H3心へ延長しない]
\label{prop:vieta-not-pointed-heart}
Casimir切断 $R_{a,b,c,D}=B_{a,b,c}/(\Omega_h+Dt^3)$ と
$R_{b,a,c,D}$ が本稿の強いH3であると仮定する。
第\ref{sec:daha-character}節の $\sigma$ は非斉次化した代数の同型を与えるが、
その同型を誘導する構造対象付きexact同値
$\qgr R_{a,b,c,D}\simeq\qgr R_{b,a,c,D}$ は存在しない。
また両Fricke中心の射影閉包が滑らかなら、中心上の $\sigma$ は
それらの間の正則同型へ延長しない。
\end{proposition}
\begin{proof}
非斉次化の次数フィルトレーションを $F_\bullet$ とする。
$\sigma(z)=(q^2-q)xy-q^2z+qc$ は正確に $F_2\setminus F_1$ に属する。
実際、Casimir関係の最高記号は正則な三次元 $xyz$ なので、
次数2以下の新しい関係を生じず、$q^2-q\ne0$ である。
一方、構造対象付きexact同値は、入力\ref{ass:recognition}(4)により
反標準環と二次包絡を復元する。
二次包絡の中央一次部分は $kt$ のみであるため、
誘導されるgraded同型は $t$ の直線を保ち、非斉次化後には $F_1$ を保つ。
これは上の $\sigma(z)$ と矛盾する。
さらに中心上でも $\sigma(Z)$ の最高記号は $e^{-1}x^3y^3$ で次数6であり、
Casimir商の先頭環でも零でない。$Z$ は次数3なので、内部自己同型を
合成してもこの矛盾は消えない。従って局所化後の自然同型の曖昧さにも依存しない。

滑らかな三次曲面では $\OO(1)=-K$ なので、射影閉包の任意の同型は
反標準埋込みの線型同型から来る。
$\sigma$ はaffine開部分間の同型を与えるので、延長があれば境界も保存し、
$\langle1,u,v,w\rangle$ を保たなければならない。
しかし $\sigma(w)=-uv-C-w$ はこの空間に属さない。
Fricke方程式は次数3であるため $uv$ を一次式へ変える関係がないからである。
\end{proof}

\begin{remark}[今後構成すべきもの]
$\sigma^2=1$ と $\rho^3=1$ を持つmodular作用と、
そのbraid群による持上げを区別する。
上の命題は非可換affine部分の対称性を否定せず、
導来自己同値や心の変更を伴う対応の非存在も主張しない。
古典的なVietaの双有理力学はCantat--Loray\footnote{S. Cantat, F. Loray,
\emph{Dynamics on Character Varieties and Malgrange irreducibility of Painlev\'{e} VI equation},
Ann. Inst. Fourier 59 (2009), 2927--2978,
\href{https://doi.org/10.5802/aif.2512}{doi:10.5802/aif.2512}, \S3.1.}
で詳しく扱われている。
本研究では、固定心上の作用を目指すのではなく、
局所化した圏、境界の双有理変更、または導来圏の中の心の変更に
対応を実現することが次の課題である。
純粋な $q$ 変形の障害を回避するには、
指定した純粋な $q$ 変形の外へ出る必要がある。
第\ref{sec:integrate}節のようにSklyanin方向を同時に動かす方法は、その具体的な候補である。
\end{remark}

\section{Casimir pencilの十個のPoisson零点とH3の二種類の$A_1$特異性}
\label{sec:casimir-poisson-leaves}
三乗根の固定包絡 $V_{0,0,0}$ では、総中心三様体の五つの節点と
Poisson括弧の零点は一致しない。この差は、同じ$A_1$特異点を持つ
三次中心切断のうちH3になるものとならないものを幾何学的に区別する。
以下のH3判定のみ、既存の固定包絡全切断分類
（\texttt{papers/sections/cubic\_quantum\_triangle\_full\_cuts.tex} の
\texttt{thm:quantum-triangle-full-H3}）に依存する。
Poisson零点とHessianの計算はその判定とは独立である。

\begin{theorem}[十個のPoisson零点と二つの五点集合]
\label{thm:ten-poisson-points}
$k$ を標数零の代数閉体、$q^2+q+1=0$、
$f=(1-q)/3$、$e=-(2q+1)/9$ とする。
$V=V_{0,0,0}$ の中心三様体は
\[
 \mathcal Y=\operatorname{Spec}
 k[X,Y,Z,\Omega]/(P),\qquad
 P=\Omega^3+f\Omega^2-XYZ+e(X^2+Y^2+Z^2).
\]
第\ref{sec:root-poisson-orders}節のPoisson括弧は
$\{X,Y\}=-9P_Z$（巡回）、$\{\Omega,-\}=0$ である。
この括弧が零となる幾何学的点の集合は、正確に次の十点である。
\begin{align*}
 p_0&=(0,0,0,0),&
 r_0&=(0,0,0,-f),\\
 p_\eta&=(2e\eta_1,2e\eta_2,2e\eta_3,-2f/3),&
 r_\eta&=(2e\eta_1,2e\eta_2,2e\eta_3,f/3),
\end{align*}
ただし$\eta_i\in\{\pm1\}$、$\eta_1\eta_2\eta_3=1$ とし、
座標順は $(X,Y,Z,\Omega)$ とする。
$p_0,p_\eta$ の五点は$\mathcal Y$の通常二重点であり、
$r_0,r_\eta$ の五点は$\mathcal Y$の\emph{滑らかな}点である。
後五点でもPoisson括弧は零であり、$\Omega:\mathcal Y^{\rm sm}\to\AA^1$
の臨界点となる。これら以外の$\mathcal Y^{\rm sm}$ではPoisson rankは2である。
ここで零点集合は幾何学的supportを意味し、零schemeの被約性は主張しない。
\end{theorem}
\begin{proof}
零点条件は$P_X=P_Y=P_Z=0$、すなわち
\[
 YZ=2eX,\qquad ZX=2eY,\qquad XY=2eZ.
\]
$e\ne0$ なので、一座標が零なら三座標とも零である。
この場合$P=\Omega^2(\Omega+f)=0$ となり、$p_0,r_0$ を得る。
三座標が非零なら三式を掛けて$XYZ=(2e)^3$、各式を用いて
$X^2=Y^2=Z^2=(2e)^2$ を得る。従って
$(X,Y,Z)=2e(\eta_1,\eta_2,\eta_3)$、$\eta_1\eta_2\eta_3=1$。
$q^2+q+1=0$ から$e^3=-f^3/27$ なので、この点で
\[
 P=\Omega^3+f\Omega^2+4e^3
   =(\Omega-f/3)(\Omega+2f/3)^2.
\]
これより残る八点を得、他に零点はない。
$P_\Omega=\Omega(3\Omega+2f)$ は$p_0,p_\eta$で零、
$r_0,r_\eta$で$f^2\ne0$ となる。
従って後五点は三様体の滑らかな点であり、前五点は特異点である。
前五点で$P$の四変数Hessian行列式はそれぞれ
$16fe^3$、$64fe^3$（各符号点で同じ）で非零である。
従って総中心の特異点は通常二重点である。
\end{proof}

\begin{theorem}[Casimir pencilにおける四つの特異中心値とH3判定]
\label{thm:casimir-four-values}
中心切断$C_D=\mathcal Y\cap\{\Omega=-D\}$を取る。
中心曲面が特異となるのは正確に
\[
 D\in\{0,\ 2f/3,\ f,\ -f/3\}
\]
であり、特異点はすべて$A_1$型で、個数とH3性は以下の通りである。
\begin{center}
\begin{tabular}{cclc}
\toprule
$D$&$A_1$の個数&総中心$\mathcal Y$上の位置&強いH3\\
\midrule
$0$&1&総中心の節点&可\\
$2f/3$&4&総中心の節点&可\\
$f$&1&総中心の滑らかな臨界点&不可\\
$-f/3$&4&総中心の滑らかな臨界点&不可\\
\bottomrule
\end{tabular}
\end{center}
最後の列は既存の固定包絡H3判定定理の適用であり、
中心のHessian計算だけで有限大域次元を結論したものではない。
特にこのCasimir pencil内では、
\emph{切断のPoisson零点が全て総中心$\mathcal Y$の特異点であること}
（零点が存在しない場合を含む）が強いH3性と同値である。
\end{theorem}
\begin{proof}
$C_D$の特異点は$P=P_X=P_Y=P_Z=0$、$\Omega=-D$ によって
与えられるから、前定理で列挙した十点のみが候補となる。
$X,Y,Z$でのHessianは
\[
 \begin{pmatrix}
 2e&-Z&-Y\\-Z&2e&-X\\-Y&-X&2e
 \end{pmatrix}.
\]
原点三座標では行列式は$8e^3\ne0$、四符号座標では
行列式は$-32e^3\ne0$ である。
従って各切断の特異点は孤立した通常二重点で、表の個数が従う。

H3の判定に用いるのは前掲の既存分類
$\Delta\ne0$かつ切断超平面$H\notin\mathcal Y^*$である。
Casimir切断では$\alpha=1,\beta=\gamma=\delta=0$、従って
$\Delta=1$ である。$r_0,r_\eta$ を通る二つの値では
$P_\Omega=f^2\ne0$ だから$H$は総中心の滑らかな点に接し、
H3ではない。
一方$p_0,p_\eta$ を通る二値では、全切断特異点は
総中心の節点に限られ、その二次切断は上記Hessianで非退化である。
従って節点での極限接超平面にも属さない。
他の滑らかな総中心点で接していれば切断に別の特異点が生じるが、
前定理によりそのような点はない。
従って$H\notin\mathcal Y^*$ であり、二つの切断は強いH3となる。
他の$D$では$C_D$は滑らかで、同様にH3判定を満たす。
\end{proof}

\begin{corollary}[一次元表現が検出する五つのPoisson零点]
\label{cor:one-dim-five-of-ten}
$V_{0,0,0}$ の一次元表現は正確に五つあり、
その中心characterは$p_0,p_\eta$ の五点である。
従って総中心の滑らかなPoisson零点$r_0,r_\eta$ は
一次元表現からは生じない。これはこれらの点上の
全ての高次元単純加群やAzumaya性を分類したという主張ではない。
\end{corollary}
\begin{proof}
一次元表現では$x,y,z$をスカラーとみなし、三関係から
$(1-q)yz=x$ と巡回二式を得る。
その全解は $(0,0,0)$ および
$(x,y,z)=(\eta_1,\eta_2,\eta_3)/(1-q)$
（$\eta_1\eta_2\eta_3=1$）である。
$X=x^3+q^2x$等と式\eqref{eq:Omega}に代入すると、中心値は
それぞれ$p_0$と$p_\eta$になる。
\end{proof}

\begin{remark}[同じADE型と異なる非可換正則性]
同じ$A_1$型の中心曲面でも、総中心の節点を切る場合は
非可換H3であり、総中心の滑らかな点に接して節点を作る場合は
H3にならない。この差は単に曲面の特異点のADE型では検出できず、
Poisson leafの総空間への埋め込み、一次元表現、局所orderを合わせて理解する必要がある。
さらに非可換表現の全fiber代数についてはBrown--Gordon型の比較が有効だが、
定理\ref{thm:casimir-four-values}だけからMorita同値・Azumaya性を
推論しない。総中心の一般係数$(a,b,c)$に対するPoisson零点の
有限被覆とそのモノドロミー、七次元の結節中心因子との関係は今後の優先課題である。
\end{remark}

\section{相対Poisson臨界スキーム、二重五次判別式、\texorpdfstring{$S_5$}{S5}モノドロミー}
\label{sec:relative-poisson-cover}
前節の十Poisson零点を、低次係数$(a,b,c)$を動かした相対スキームとして調べる。
以下では被約supportとスキームの長さを区別する。基礎体$k$は標数0の代数閉体、
$q^2+q+1=0$、$f=(1-q)/3$、$e=-(2q+1)/9$とする。

\subsection{全係数で有限平坦な二つの被覆}
$A=k[a,b,c]$、$S=\Spec A$とし、式\eqref{eq:center}の左辺を
$P_{\mathbf a}(X,Y,Z,\Omega)$と書く。Casimir切断の特異点を添字付ける
相対Poisson臨界スキームは
\[
 \mathcal K=\Spec\frac{A[X,Y,Z,\Omega]}
 {(P_{\mathbf a},P_X,P_Y,P_Z)}\xrightarrow{\kappa}S
\]
である。ここで偏微分は$X,Y,Z$に関するもので、$\Omega$は固定しない。

\begin{theorem}[次数十五と五の全係数有限平坦性]
\label{thm:relative-critical-flat}
$\kappa$は有限平坦次数$15$である。また一次元表現の相対スキーム
\[
 \mathcal E=\Spec\frac{A[x,y,z]}
 {((1-q)yz-x-a,\ (1-q)zx-y-b,\ (1-q)xy-z-c)}
 \xrightarrow{\pi}S
\]
は有限平坦次数$5$であり、総空間は$\AA^3_{x,y,z}$と同型なので既約である。
$\pi$は原点を含む非空Zariski開集合上で有限\'{e}taleとなる。
\end{theorem}
\begin{proof}
$P_{\mathbf a}$を通常の全次数で斉次化し、
$[X:Y:Z:\Omega:T]\in\PP^4_S$で三次式$P^h$と
三二次式$P_X^h,P_Y^h,P_Z^h$の共通零点を$\overline{\mathcal K}$とする。
$T=0$では
\[
 (P^h,P_X^h,P_Y^h,P_Z^h)
 =(\Omega^3-XYZ,-YZ,-ZX,-XY),
\]
従ってsupportは三座標頂点のみである。
$X=1$の頂点近傍では
\begin{align*}
P_X^h&=-YZ+T(2e-a\Omega)+T^2\ell_a,\\
P_Y^h&=-Z+T(2eY-b\Omega)+T^2\ell_b,\\
P_Z^h&=-Y+T(2eZ-c\Omega)+T^2\ell_c .
\end{align*}
これらの$(T,Z,Y)$に関するJacobianは頂点で可逆（行列式$\pm2e$）。
$T=Y=Z=0$が$\Omega$の任意値に対する厳密解なので形式陰関数定理から、
三式で割った完備局所環では$T=Y=Z=0$、残る関係は$\Omega^3=0$となる。
よって各無限遠頂点は相対長さ3の自明な厚みである。

任意の正次元射影部分多様体は無限遠と交わるが、上の三頂点は孤立している。
従って全幾何学的fiberは零次元で、$\overline{\mathcal K}$は
次数$(3,2,2,2)$の相対完全交叉である。
B\'{e}zoutから各fiberの長さは$3\cdot2^3=24$。
射影かつ準有限だから有限であり、完全交叉のCohen--Macaulay性と
regularな底$S$へのmiracle flatnessにより平坦である。
無限遠の三つの厚みは開閉な有限平坦部分（次数9）なので、
その補集合である$\mathcal K$は有限平坦次数$24-9=15$となる。

一次元表現について三二次関係を$[x:y:z:T]\in\PP^3_S$に斉次化する。
$T=0$では三座標頂点のみで、例えば$x=1$で
\[
 (1-q)yz-T-aT^2,\quad (1-q)z-yT-bT^2,\quad
 (1-q)y-zT-cT^2
\]
の$(T,z,y)$に関するJacobianは単元$\pm(1-q)^2$である。
三つの無限遠点は長さ1の自明な開閉section。
射影三二次完全交叉は有限平坦次数$8$なので、これらを除く
$\mathcal E$は次数$8-3=5$となる。
$a,b,c$は$x,y,z$の多項式だから総空間は$\AA^3$。
原点fiberの五点のJacobianは$-1,4,4,4,4$で非零
（原点と $(x,y,z)=(\eta_1,\eta_2,\eta_3)/(1-q)$、$\eta_i=\pm1$、$\eta_1\eta_2\eta_3=1$ への直接代入）なので\'{e}tale開集合は非空である。
\end{proof}

\subsection{二重五次因子とコンパニオン五次因子}
各一次元表現に$X,Y,Z,\Omega$を評価すると、
全係数で$P=P_X=P_Y=P_Z=P_\Omega=0$が成立する。
これは元関係の代入恒等式であり、
\texttt{computations/verify\_character\_node\_cover.py}でも検証済みである。

\begin{theorem}[二重節点とCasimir臨界値の五次因数分解]
\label{thm:critical-discriminant-quintics}
原点を含む非空Zariski開近傍$S^\circ$上で、$\mathcal K$は
相対節点に沿う次数10の二重厚みと、総中心の滑らかなPoisson臨界点を
添字付ける有限\'{e}tale五重被覆$\mathcal E'\to S^\circ$に開閉分解する。
前者は$\mathcal E|_{S^\circ}$の相対一次無限小厚みで、
\'{e}tale局所的には
\[
 \mathcal K|_{S^\circ}\simeq
 \Spec_{\mathcal E}(\OO_{\mathcal E}[\epsilon]/\epsilon^2)
 \sqcup\mathcal E'.
\]
$\mathcal E'$の幾何学的既約性・monodromyは、後述の定理\ref{thm:paired-quintic-S5}で完全決定する。

二つのmonicな特性多項式
\[
 N_5(D)=\det(D+\Omega\mid\pi_*\OO_{\mathcal E}),\qquad
 \Delta_{15}(D)=\det(D+\Omega\mid\kappa_*\OO_{\mathcal K})
\]
は$A[D]$上の次数5と15の多項式であり、
\begin{equation}\label{eq:relative-discriminant-quintics}
 \boxed{\Delta_{15}(D)=N_5(D)^2M_5(D)},\qquad
 M_5\in A[D]\text{ monic of degree }5
\end{equation}
が\emph{全係数}で成立する。
$\Delta_{15}=0$はCasimir中心切断が特異となる値の条件、
$N_5=0$は一次元表現に由来する節点の中心値である。
原点での特殊化は
\[
 N_5(D)=D(D-2f/3)^4,\qquad
 M_5(D)=(D-f)(D+f/3)^4.
\]
\end{theorem}
\begin{proof}
原点fiberの臨界点supportは定理\ref{thm:ten-poisson-points}の十点。
全点で$(X,Y,Z)$-Hessianが非退化なので、$P_X=P_Y=P_Z=0$から
$X,Y,Z$を$(\Omega,a,b,c)$の形式関数に解ける。
これを$P$に代入した一変数方程式を$h(\Omega)$とする。
総中心の五節点では$P_\Omega=0$、$P_{\Omega\Omega}=2f,-2f$は単元。
一次元表現による節点sectionは全係数で
$P=P_X=P_Y=P_Z=P_\Omega=0$を満たして存続する。
従って各節点上で$h=(\Omega-\Omega_{\rm node})^2\cdot\mathrm{unit}$となる。
他の五点では$P_\Omega=f^2$が単元なので、四式のJacobianが可逆で、
相対長さ1の有限\'{e}tale点へ延長する。
有限平坦性と原点の十点の相異性から、底を開に縮めればこれらは
互いに素な開閉部分へ分離する。次数は$5\cdot2+5=15$である。

この分解における$D+\Omega$の特性多項式は
相対節点側で$N_5(D)^2$、残部でmonic五次式になる。
従って$S^\circ$上で式\eqref{eq:relative-discriminant-quintics}が成立する。
$A=k[a,b,c]$は整域で$N_5$はmonicなので、
$A[D]$で$\Delta_{15}$を$N_5^2$で割った剰余が
稠密開上で零なら全域で零となる。ゆえに$M_5\in A[D]$が一意に存在する。
原点での式は、五つの$(X,Y,Z)$-臨界三つ組に対する$\Omega$多項式が
一つは$\Omega^2(\Omega+f)$、残る四つは
$(\Omega-f/3)(\Omega+2f/3)^2$であることから従う。
\texttt{verify\_relative\_poisson\_cover.py}は
$\mathbb Q(\sqrt{-3})$上のGr\"obner商の長さ15、
$-\Omega$の$15\times15$乗法行列の特性多項式、
無限遠のJacobianを独立に厳密検算する。
\end{proof}

\subsection{既約五次因子の\texorpdfstring{$S_5$}{S5}モノドロミー}
\begin{theorem}[一次元表現被覆の完全対称群]
\label{thm:character-cover-S5}
$\mathcal E\to S$の最大\'{e}tale開部分の幾何学的monodromy群は
五文字の完全対称群$S_5$である。
さらに$N_5(D)$は$k(a,b,c)[D]$上既約・分離的で、
その分解体のGalois群も$S_5$である。
\end{theorem}
\begin{proof}
$r=1-q$、$u=rx,v=ry,w=rz$、$A=ra,B=rb,C=rc$と置くと、
一次元表現の係数写像は
\begin{equation}\label{eq:five-character-gradient}
 (u,v,w)\mapsto(A,B,C)=(vw-u,wu-v,uv-w)
\end{equation}
となる。これは$uvw-(u^2+v^2+w^2)/2$のgradient mapである。
被覆の総空間は既約な$\AA^3_{u,v,w}$なのでmonodromyは五点に推移的であり、
次数が素数5なのでprimitiveである。
Jacobian行列と行列式は
\[
 J=\begin{pmatrix}-1&w&v\\w&-1&u\\v&u&-1\end{pmatrix},
 \quad\det J=u^2+v^2+w^2+2uvw-1.
\]
$A=B=C=-1/4$のfiberには$(1/2,1/2,1/2)$が重複度2で現れ、
他は$(5/4,-1,-1)$の三順列（Jacobian行列式$81/16$）である。
前者で$J$の核・余核は対角方向で、対角写像
$t\mapsto t^2-t$は$t=1/2$で単純foldとなる。
従って幾何学的慣性群には一つの互換がある。
primitive置換群において一つの互換の全共役を辺とするグラフは
連結でなければ非自明なblock系を作る。
連結グラフの辺に沿う互換は$S_5$を生成するから、
monodromy群は$S_5$である。

一次元表現条件を式\eqref{eq:Omega}へ代入し
$q^2+q+1=0$で還元すると、
\begin{equation}\label{eq:five-character-casimir}
 \Omega_{\mathcal E}(u,v,w)=\frac f3(uvw-u^2-v^2-w^2)
\end{equation}
を得る。原点fiberの五点は$(0,0,0)$および
$\eta=(\eta_1,\eta_2,\eta_3)$、$\eta_i=\pm1$、積1である。
前者の$\Omega$値は$0$、後四点は$-2f/3$だが、
その底$(A,B,C)$上の一次微分は
\[
 d\Omega_\eta=-\frac f3(\eta_1dA+\eta_2dB+\eta_3dC).
\]
これは$\nabla_{u,v,w}\Omega_\eta=-f\eta/3$かつ$J(\eta)\eta=\eta$から従う。
四つの一次形式は相異なるので、五つの$\Omega$値は
一般の底上で相異なる。従って$\Omega_{\mathcal E}$は
次数5の拡大$k(\mathcal E)/k(a,b,c)$の原始元で、
$N_5$は既約・分離的、そのGalois群は先ほどの$S_5$である。
\texttt{verify\_relative\_poisson\_cover.py}はsimple fold、
Casimirの式と一次微分を全係数で厳密検算する。
\end{proof}

\subsection{二つの五次因子の独立な\texorpdfstring{$S_5$}{S5}モノドロミー}
\begin{theorem}[コンパニオン五次式の幾何学的$S_5$と独立性]
\label{thm:paired-quintic-S5}
式\eqref{eq:relative-discriminant-quintics}の二つの五次因子
$N_5(D),M_5(D)$について、標数零の任意の代数閉体$k$上で次が成り立つ。
\begin{enumerate}
\item $M_5(D)$は$k(a,b,c)[D]$上既約かつ分離的であり、
その幾何学的Galois群は$S_5$である。
従って定理\ref{thm:critical-discriminant-quintics}の
残余有限\'{e}tale五重被覆$\mathcal E'$は幾何学的に連結である。
\item $N_5,M_5$の分解体をそれぞれ$L_N,L_M$とすると
\[
 \operatorname{Gal}(L_NL_M/k(a,b,c))\cong S_5\times S_5.
\]
特に両分解体は線型非交叉で、二つの判別式が定める
$k(a,b,c)^\times/k(a,b,c)^{\times2}$の平方類も独立である。
\end{enumerate}
この結果は二つの被覆のGalois幾何を決めるが、
一般$(a,b,c)$の非可換H3性・全Azumaya locus・局所Morita型を
それだけで判定するものではない。
\end{theorem}
\begin{proof}
まず$K_0=\mathbb Q(q)$で考える。
全係数の有限平坦性は$\mathbb Z[q,1/6]$上でも成立し、
$q^2+q+1=0$を用いた係数表示の分母は$3$の冪なので、
以下の素数における有限体還元はよく定義される。
$K_0$-有理な係数$(a,b,c)=(30,19,29)$で特殊化する。
次の二つの根の選択は、それぞれ$K_0$の素イデアル
$q\equiv374\pmod{1009}$、$q\equiv86\pmod{1069}$を指定する。
両方で$q^2+q+1\equiv0$が成立する。

各有限体上で$P,P_X,P_Y,P_Z$のGr\"obner商の長さが15であることを
確認し、$-\Omega$乗法の$15\times15$行列を構成する。
これとは独立に、一次元表現の三二次関係のGr\"obner商の長さ5を確認し、
式\eqref{eq:five-character-casimir}の$-\Omega$乗法行列から
$N_5(D)$を得る。前者の特性多項式を$N_5(D)^2$で割ると
余り零、商は$M_5(D)$となる。正確な結果は次の通りである。
\begin{align*}
 M_5\bmod(1009,q-374)
  ={}&D^5-404D^4-215D^3+279D^2+226D+42,\\
 M_5\bmod(1069,q-86)
  ={}&(D^2+248D+337)(D^3+55D^2-276D-260).
\end{align*}
前者は$\mathbb F_{1009}$上既約、後者の二因子は
$\mathbb F_{1069}$上それぞれ既約で互いに異なる。
後者の素数では$N_5$は
\[
 D^5-100D^4-59D^3-17D^2+487D+534
\]
であり、こちらも$\mathbb F_{1069}$上既約で、$M_5$と互いに素である。
両素数で多項式の判別式は零でない。
これらの有限体上の等式、因数分解、既約性は
\texttt{computations/verify\_companion\_s5.py}で
元の中心と一次元表現の関係式から正確に再生成できる。

有限体のFrobeniusの分解型はそれぞれ$(5)$と$(2,3)$なので、
$K_0$上の特殊化$M_5$の算術Galois群には
位数5の元と型$(2,3)$の元がある。
後者は奇置換かつ位数6である。
従って群位数は$30$の倍数で、$S_5$の部分群だから
$30,60,120$のいずれかである。
位数60なら指数2の部分群$A_5$となり奇置換を含めない。
位数30なら$S_5$の指数4の部分群を与えるが、
四剰余類への推移的作用$S_5\to S_4$は存在しない
（$S_5$の固有非自明正規部分群は$A_5$のみ）。
よってGalois群は$S_5$である。
これを一般係数へ特殊化のGalois群の包含により戻すと、
$K_0(a,b,c)$上の$M_5$の算術Galois群も$S_5$である。

ここで算術的結論と幾何学的結論を区別する。
定理\ref{thm:critical-discriminant-quintics}の
$K_0$上の非空Zariski開集合$S^\circ$は係数原点を含み、
$\mathcal E'\to S^\circ$は有限\'{e}tale次数5である。
原点の五点は定理\ref{thm:ten-poisson-points}の
$r_0,r_\eta$で、すべて$K_0$-有理かつ相異なる。
従って$\mathcal E'$の原点fiberは$K_0^5$に完全分裂する。
$S^\circ$は幾何学的に連結であるから、\cite[Lemma 58.14.3]{StacksPi1}
の算術・幾何基本群完全列
\[
 1\longrightarrow\pi_1(S^\circ_{\overline K_0})
 \longrightarrow\pi_1(S^\circ)
 \longrightarrow\operatorname{Gal}(\overline K_0/K_0)\longrightarrow1
\]
を使う。原点という$K_0$-有理基点は右側へのsectionを与え、
$\mathcal E'$の五点fiberが完全分裂しているので
そのsectionの$S_5$への作用は自明である。
従って算術的な$S_5$作用は幾何基本群の像そのものであり、
$M_5$の幾何学的Galois群も$S_5$となる。
ここで$M_5(D)$自体は原点で四重根を持つが、
\emph{エタールな被覆は多項式の零点スキームではなく
$(X,Y,Z,\Omega)$を保持した$\mathcal E'$である}点に注意する。
これが(1)を示す。

(2)を示す。$N_5$側の算術・幾何群は
定理\ref{thm:character-cover-S5}により$S_5$であり、
$M_5$側も上で確定した。合併した算術monodromy
$G\subset S_5\times S_5$は両射影へ全射である。
$p=1069$では$N_5$の分解型が$(5)$で偶置換、
$M_5$の分解型が$(2,3)$で奇置換なので、
$G$には符号の異なる対が存在する。
Goursatの補題と$S_5$の正規部分群
$\{1\},A_5,S_5$から、真のsubdirect productは
二つの符号が一致する$C_2$上のfiber product、または
$S_5$の自己同型のグラフに限られる。
いずれも符号の異なる対を含まないから
$G=S_5\times S_5$である。

幾何学的同時monodromyについては、同じ有理基点における
$\mathcal E$の五点も$p_0,p_\eta$として全て$K_0$-有理である。
よって$\mathcal E\sqcup\mathcal E'$の十点fiberは完全分裂し、
基本群完全列の有理sectionは$S_5\times S_5$への作用が自明。
従って算術と幾何の同時monodromyは一致する。
特に合成分解体は$S_5\times S_5$をGalois群に持ち、
二分解体は線型非交叉である。
その二つの符号商$S_5\to C_2$は互いに独立だから、
判別式の平方類も独立となる。
$\overline K_0$から任意の標数零代数閉体$k$への拡大でも、
有限型Galois閉包の幾何学的連結性は保たれる。
\end{proof}

\begin{remark}[計算の意味と未解決事項]
二つの$S_5$の独立性は有限体の無作為標本からの外挿ではない。
全係数で定義したmonic五次式について、二つの指定した良い素数における
\emph{厳密な}因数分解からFrobeniusの置換型を証明し、
群論とエタール基本群の完全列を適用したものである。
これによって残余五次式$M_5$の幾何学的既約性・monodromyは解決した。
未決定なのは両五次式の具体的な大域係数の簡潔な表示、
分岐divisorの幾何学、局所order/H3判定との一般係数での関係、
および楕円$E_6$成分への移行である。
\end{remark}

\begin{remark}[Painlev\'{e} VIの判別式との区別]
Inaba--Iwasaki--Saito \cite[\S9.1, Definition 9.1]{IISDiscriminant}は
Fricke cubicの判別式を境界トレースの係数空間に引き戻すと
二乗因子が現れることを示す。
本稿の$N_5^2M_5$は\emph{別のパラメータ写像}の
次数十五の臨界値多項式であり、両者の因子が等しいとは主張しない。
五次式$M_5$の既約性・幾何学的$S_5$および$N_5$との独立性は
定理\ref{thm:paired-quintic-S5}で解決した。
一般$(a,b,c)$でのH3条件、局所order、楕円境界の八次元成分への
相対中心の延長と分岐divisorの比較は未解決である。
本稿の計算は既知のPainlev\'{e}判別式の再発見とは区別し、
個別の新規性は一次文献との詳細照合を待つ。
\end{remark}

\section{五重被覆の内側に現れるCayley曲面、$D_4$特異点、実凸幾何}
\label{sec:cayley-lg-connections}
本節は、一次元表現の五重被覆を異なる理論の共通の具体例として読む。
対応の出発点は、正規化した係数と一次元表現
\[
 (A,B,C)=(1-q)(a,b,c),\qquad (u,v,w)=(1-q)(x,y,z)
\]
に関する一つの写像である：
\begin{equation}\label{eq:cayley-gradient-map}
 \Phi(u,v,w)=(vw-u,wu-v,uv-w)=\nabla H,
 \qquad H=uvw-\frac12(u^2+v^2+w^2).
\end{equation}
本節の$A,B,C$はこの\emph{正規化した低次係数}であり、
第\ref{sec:daha-character}節のFricke三係数とは異なる。
また$u,v,w$は中心生成元$X,Y,Z$ではない。
以下で扱う三つの空間を区別する：
一次元表現の総空間$\AA^3_{u,v,w}$、係数底$\AA^3_{A,B,C}$、
および元の非可換代数の三次中心。

\subsection{Landau--GinzburgポテンシャルとLagrangian射影}
\begin{proposition}[一次元表現は五つの臨界点を与える]
\label{prop:lg-character-potential}
\[
 W_{A,B,C}(u,v,w)=H(u,v,w)-Au-Bv-Cw
\]
と置く。一次元表現の相対スキームは、この多項式族の相対臨界スキームであり、
\[
 \operatorname{Jac}(W)=
 k[A,B,C,u,v,w]/(vw-u-A,wu-v-B,uv-w-C)
\]
は底上有限平坦rank $5$である。
また$T^*\AA^3_{u,v,w}$内の$\operatorname{graph}(dH)$はexact Lagrangianで、
余接座標$(A,B,C)$への射影が$\Phi$に一致する。
\end{proposition}
\begin{proof}
偏微分を取ると三つの一次元表現方程式がそのまま出る。
rank $5$の有限平坦性は定理\ref{thm:relative-critical-flat}による。
Liouville形式$\lambda=A\,du+B\,dv+C\,dw$のgraphへの制限は$dH$であり、
従って$d\lambda$の制限は零である。graphの次元は半次元の$3$なので結論を得る。
\end{proof}

\subsection{分岐源はトーラスの指標多様体そのものである}
\begin{theorem}[Cayley曲面とトーラス商]
\label{thm:ramification-cayley-torus}
$\Phi$のramification divisorは
\begin{equation}\label{eq:cayley-ramification}
 \mathcal R=\{R=0\},\qquad
 R=u^2+v^2+w^2+2uvw-1.
\end{equation}
$x=-2u,y=-2v,z=-2w$により、これはCayley三次曲面
\[
 x^2+y^2+z^2-xyz-4=0
\]
となり、粗指標多様体
$\operatorname{Hom}(\mathbb Z^2,\mathrm{SL}_2)\mathbin{/\!/}\mathrm{SL}_2$
と同型である。具体的には
\begin{equation}\label{eq:cayley-torus-cover}
 (r,t)\longmapsto
 \left(-\frac{r+r^{-1}}2,-\frac{t+t^{-1}}2,
 -\frac{rt+(rt)^{-1}}2\right)
\end{equation}
が
\[
 (\mathbb G_m)^2/\langle(r,t)\mapsto(r^{-1},t^{-1})\rangle
 \xrightarrow{\sim}\mathcal R
\]
を与える。$\mathcal R$の四つの特異点は
$\sigma_i\in\{\pm1\}$、$\sigma_1\sigma_2\sigma_3=-1$の点で、全て$A_1$である。
\end{theorem}
\begin{proof}
\[
 J=\operatorname{Hess}H=
 \begin{pmatrix}-1&w&v\\w&-1&u\\v&u&-1\end{pmatrix},\qquad \det J=R
\]
から分岐源の式を得る。変数変換によるCayley方程式は直接展開で確かめられる。
トーラス上の同時反転不変環は
$r+r^{-1},t+t^{-1},rt+(rt)^{-1}$で生成される。
実際、対称Laurent単項式
$L_{m,n}=r^mt^n+r^{-m}t^{-n}$に対して
$L_{1,0}L_{m,n}=L_{m+1,n}+L_{m-1,n}$と巡回的な漸化式を使い、
$L_{1,-1}=L_{1,0}L_{0,1}-L_{1,1}$から全てを生成できる。
唯一の関係は表示した三次式である。
可換な半単純$\mathrm{SL}_2$表現を対角化すると、残余Weyl作用が同時反転となる。
従ってこの粗商が閉トーラスの指標多様体である。
同時反転の固定点$(r,t)\in\{\pm1\}^2$の局所商は
$\AA^2/(\pm1)$なので四つの$A_1$を与える。
それ以外では商写像は\'{e}taleである。
\end{proof}

このCayley曲面のトーラス商表示とmodular力学自体は既知である
\cite[\S1.5]{CantatLorayCayley}。
本稿での接続は、\emph{非可換三次曲面の一次元表現を添字付ける
五重被覆の分岐源に、この別の指標多様体が現れる}という点にある。
$\mathcal R$と、その像であるbranch locus
$\mathcal B=\Phi(\mathcal R)\subset\AA^3_{A,B,C}$を同一視しない。

\begin{proposition}[対数Poisson構造とmodular対称性]
\label{prop:ramification-log-poisson}
トーラスに$\{r,t\}=rt$を入れると、式\eqref{eq:cayley-torus-cover}の座標は
\[
 \{u,v\}=-\frac12 R_w,\qquad
 \{v,w\}=-\frac12 R_u,\qquad
 \{w,u\}=-\frac12 R_v
\]
を満たす。$\mathrm{GL}_2(\mathbb Z)$の単項式作用は同時反転と可換し、
$\mathrm{PGL}_2(\mathbb Z)$の作用として$\mathcal R$に降りる。
行列式$1$の元はPoisson、行列式$-1$の元は反Poissonである。
\end{proposition}
\begin{proof}
例えば$\{u,v\}=(r-r^{-1})(t-t^{-1})/4$であり、
$R_w=2w+2uv=-(r-r^{-1})(t-t^{-1})/2$となる。
他の二式も直接計算による。単項式変換の対数Jacobianはその整数行列なので、
$d\log r\wedge d\log t$は行列式倍される。
$-I$は同時反転そのものであるから商への作用は$\mathrm{PGL}_2$を経由する。
\end{proof}
ここでの括弧は\emph{分岐源}上のもので、元の量子代数の\emph{中心}上の
root-of-unity括弧の係数$9$とは異なる空間・規約である。
またこのmodular作用が五重被覆全体と係数底への作用へ延長するとは主張しない。

\subsection{四枚の衝突は$D_4$ミニバーサル変形である}
\begin{theorem}[四つの$D_4$点と一つの非分岐sheet]
\label{thm:character-D4-unfolding}
$\mathcal R$の四特異点$\sigma$の像は$\Phi(\sigma)=-2\sigma$である。
各fiber $\Phi^{-1}(-2\sigma)$は、$\sigma$に支持された長さ$4$の局所スキームと、
点$-2\sigma$に支持された一つの単純点からなる。
前者で$W_{A,B,C}$は$D_4$型の孤立臨界点を持ち、三つの係数変形は
加法定数を除いたミニバーサル変形を与える。
加法定数も独立に加えれば、通常の四変数ミニバーサル底を得る。
\end{theorem}
\begin{proof}
二符号変更により、$\sigma=(-1,-1,-1)$、$(A,B,C)=(2,2,2)$の場合で十分である。
$u=-1+x,v=-1+y,w=-1+\ell-x-y$、
$A=2+\alpha,B=2+\beta,C=2+\gamma$と置く。
さらに$L=\ell-xy+\gamma$とすると、\emph{正確な多項式恒等式}
\begin{align}\label{eq:D4-exact-splitting}
 W={}&-\frac12L^2-xy(x+y)+\frac12x^2y^2
       -\gamma xy-(\alpha-\gamma)x-(\beta-\gamma)y\notag\\
    &+\frac72+\alpha+\beta+\gamma+\frac12\gamma^2
\end{align}
を得る。
残った二変数germ $g=-xy(x+y)+x^2y^2/2$の三次部分は
相異なる三直線の積である。Morseの分離補題と$D_4$の標準的分類により、
これは$D_4$の一変数suspensionである。
例えば三次部分のJacobi環では
$x^2=y^2=-2xy$、$x^2y=xy^2=0$で、基底は$1,x,y,xy$となる。
四次項を含む$g$の原点局所Jacobi環でも同じ基底と関係が成立する。
この事実は二偏微分の式から求めることも、他の臨界点を除いた
局所化により直接計算することもできる。
従ってMilnor数は$4$であり、$D_4$の$3$-決定性により四次項は右同値で除ける。

式\eqref{eq:D4-exact-splitting}の係数微分は、定数を除いて
$-x,-y,x+y-xy$となり、局所Jacobi環の$\langle x,y,xy\rangle$を張る。
無限小versal性判定から、三係数は加法定数を除くミニバーサル変形となる。
この判定と$D_4$の標準的変形については
\cite[Proposition 3, \S5]{MikoszWeberTriality}も参照されたい。

同fiberの別の点$(2,2,2)$では$\det J=27\ne0$で、長さ1の単純点である。
既存の全域次数5の有限平坦性と$4+1=5$から、他の点は存在しない。
他の三つの$\sigma$は二符号変更で得る。
\end{proof}

\begin{corollary}[trialityの具体的な入口]
\label{cor:gradient-local-triality}
$u,v,w$の巡回置換は、$(A,B,C)=(2,2,2)$の$D_4$ germを保存する。
Hessian核$\ell=0$上では
\[
 (x,y)\longmapsto(y,-x-y)
\]
として作用し、三次接錐の三枝を巡回する。この作用は位数3である。
\end{corollary}
\begin{proof}
$H$および係数対角上の$Au+Bv+Cw$は巡回置換不変である。
核平面の座標は$(x,y,-x-y)$なので表示した線型作用を得る。
$-xy(x+y)$を保存し三枝を巡回することは直接確かめられる。
\end{proof}
これは$D_4$ trialityとその$G_2$型折り畳みを調べる明示的な足掛かりであり、
その一般理論は\cite{MikoszWeberTriality}にある。
ただし本稿で以前得た三次曲面のPicard格子内の$D_4$と、
この局所Milnor格子との\emph{標準的な同定}はまだ構成していない。
Lie群、vanishing cycle、非可換圏の作用を名称の一致から同一視しない。

\subsection{Jacobi代数、留数対、判別式を読む際の区別}
\begin{proposition}[原点のFrobenius代数と非分岐な値の衝突]
\label{prop:jacobi-residue-maxwell}
係数原点では$\Phi^{-1}(0)$は
$0$と$\eta_i\in\{\pm1\}$、$\eta_1\eta_2\eta_3=1$の五点からなる。
このJacobi代数に
\[
 \tau(h)=-h(0)+\frac14\sum_\eta h(\eta),\qquad
 (h_1,h_2)=\tau(h_1h_2)
\]
を入れると、対称非退化なFrobenius代数となる。
これは非退化臨界点における$1/\det\operatorname{Hess}W$重みの留数対である。
一方、同じ原点では五重被覆は\'{e}taleであるにもかかわらず、
Casimir normは$N_5(D)=D(D-2f/3)^4$となる。
従って$\operatorname{Disc}_D N_5=0$を被覆$\Phi$の分岐条件と同一視できない。
\end{proposition}
\begin{proof}
五点でHessian行列式は$-1,4,4,4,4$である。
Jacobi代数は五点上の関数環$k^5$なので、全て非零の重みを持つ上の対は非退化。
基底$1,u,v,w,uvw$では対の行列式は$-1$となる。
Casimirの五点での値は$0,-2f/3,-2f/3,-2f/3,-2f/3$であり、後半も従う。
\end{proof}

この差は二つの判別式を比較する際に重要である。
一般に有限局所自由代数$B/R$の基底$e_1,\ldots,e_n$と元$\omega$を取り、
$(1,\omega,\ldots,\omega^{n-1})=(e_1,\ldots,e_n)T$と書けば、
trace pairingの行列式について
\[
 \operatorname{Disc}(1,\omega,\ldots,\omega^{n-1})
     =(\det T)^2\operatorname{Disc}(e_1,\ldots,e_n)
\]
が成り立つ。これは行列$T^{\mathsf t}GT$の行列式を取った恒等式である。
Casimirという一つの観測量がsheetを区別しなくなる因子$\det T=0$は、
被覆自体の分岐とは別に現れる。
Landau--Ginzburgの臨界値が衝突するMaxwell locusとの比較も有望だが、
本稿のCasimir値は一般には$W$の臨界値ではない。
実際、臨界点上では
\[
 W_{A,B,C}=-2uvw+\frac12(u^2+v^2+w^2),\qquad
 \Omega_{\mathcal E}=\frac f3(uvw-u^2-v^2-w^2)
\]
であり、同じ関数として扱えない。
Frobenius\emph{代数}の構成だけから、primitive form、Frobenius\emph{多様体}、
またはFukaya--Seidel圏との同値を結論しない。

\subsection{実型はelliptopeとpillowcaseを結ぶ}
正規化した実多項式$\Phi:\mathbb R^3\to\mathbb R^3$を考える。
これは$q$自体が実であるという意味ではない。
\begin{proposition}[相関行列の凸領域と一意な実sheet]
\label{prop:elliptope-gradient-injective}
\[
 \mathcal E_3^-=
 \left\{(u,v,w)\in\mathbb R^3:
 C(u,v,w)=\begin{pmatrix}1&-w&-v\\-w&1&-u\\-v&-u&1\end{pmatrix}\succeq0\right\}
\]
と置く。これは符号・座標順を変えた$3$-elliptopeで、
\[
 \det C=-R,\qquad
 \partial\mathcal E_3^-=\mathcal R(\mathbb R)\cap[-1,1]^3.
\]
その境界はトーラス表示\eqref{eq:cayley-torus-cover}に
$|r|=|t|=1$を入れた像であり、
$\mathrm{SU}(2)$トーラス指標多様体のpillowcaseとなる。
$\mathcal E_3^-$の内部では$\Phi$は単射である。
従って任意の実係数に対し、この領域内の実一次元表現は高々一つである。
\end{proposition}
\begin{proof}
半正定値行列の凸性から$\mathcal E_3^-$は凸である。
主小行列式条件は$|u|,|v|,|w|\le1$と$-R\ge0$に等しい。
正定値行列が内部であり、その境界はdeterminant零だから表示を得る。
境界点は三つの単位ベクトルの階数$2$以下のGram行列である。
それらを平面上に配置し、角度を$0,\theta+\varphi,\varphi$と取ると
$(u,v,w)=(-\cos\theta,-\cos\varphi,-\cos(\theta+\varphi))$となる。
逆も明らかで、$(S^1)^2$の同時反転商がこの境界に一致する。

内部では$\operatorname{Hess}H=-C$が負定値である。
異なる二点$p,q$を結ぶ線分も内部にあり、
\[
 (\Phi(p)-\Phi(q))\cdot(p-q)
 =\int_0^1(p-q)^{\mathsf t}\operatorname{Hess}H(q+t(p-q))(p-q)\,dt<0.
\]
従って$\Phi(p)\ne\Phi(q)$となる。
\end{proof}

このelliptopeとCayley曲面は相関行列・半正定値最適化の基本例であり、
Dixon \cite[Lemma 3.5]{DixonElliptope}は同じ凸体を
nearly K\"ahler $S^3\times S^3$のmulti-moment mapの像として同定している。
ここで得た共通の凸体の同定は、非可換曲面とnearly K\"ahler多様体の
圏同値やmetric degenerationの構成ではない。
一方、実sheetの一意性という具体的帰結を与え、
実表現・安定性・pillowcaseの幾何を比較する入口にはなる。

\begin{remark}[確立した接続と研究の射程]
Cayley曲面のトーラス商、$D_4$の標準変形、elliptopeと
multi-moment像の関係は、それぞれの分野で既知である。
本節はそれらを\emph{同一の一次元表現被覆の具体式}を介して結び付け、
$D_4$の四sheet衝突・係数のversal性・実sheet一意性を証明した。
各個別式の文献上の新規性は未確認である。
既存の$S_5\times S_5$定理やH3分類を、新しい曲面・特異点との
同型だけから強化するものではない。
次に構成すべき対象は、局所vanishing cycleとPicard格子の比較、
trialityの圏的実現、留数対の変形、および分岐とCasimir値衝突の分離である。
\end{remark}

\section{消滅サイクルからPicard格子へ、そして例外群\texorpdfstring{$G_2$}{G2}へ}
\label{sec:lg-picard-g2}
本節は$k=\mathbb C$とし、
\[
 W_{A,B,C}=uvw-\tfrac12(u^2+v^2+w^2)-Au-Bv-Cw,\qquad
 U_{A,B,C;h}=\{W_{A,B,C}=h\}
\]
を考える。前節の一次元表現の\emph{臨界点被覆}と、本節の
\emph{二次元の等位面}を区別する。
Goldman--NeumannによるFricke曲面の位相、Deligneの混合Hodge理論、
Boalch--Palubaの$G_2$指標多様体の定理を用い、
この模型に現れる二つの$D_4$を実際の写像で結ぶ。
一般理論自体は既知であり、以下の役割はその適用と正規化の明示である。

\subsection{LG等位面はFricke族全体を与える}
\begin{proposition}[正規化と三角境界]
\label{prop:lg-fricke-compactification}
$p=-2u,\ r=-2v,\ s=-2w$と置けば、$U_{A,B,C;h}$は
\begin{equation}\label{eq:lg-fricke-level}
 p^2+r^2+s^2+prs-4Ap-4Br-4Cs+8h=0
\end{equation}
となる。射影閉包$S\subset\PP^3$は
\[
 uvw-\tfrac12t(u^2+v^2+w^2)-t^2(Au+Bv+Cw)-ht^3=0
\]
であり、無限遠$D=S\cap\{t=0\}=D_1+D_2+D_3$は三角形である。
$S$は$D$に沿って常に滑らかで、$U$が滑らかなら$S$は滑らかな三次曲面となる。
\end{proposition}
\begin{proof}
式\eqref{eq:lg-fricke-level}の左辺に$(-2u,-2v,-2w)$を代入すると
$-8(W-h)$となる。無限遠の辺の内部では反対の座標による偏微分が非零、
各頂点では$t$微分が非零の平方の$-1/2$倍となる。
\end{proof}

元の量子代数のCasimir切断との比較も具体的である。
その中心の記号$f,e,\ell_a,\ell_b,\ell_c,K,L,D_0$を用い、
\[
 A=\frac{D_0a+\ell_a}{4e^2},\quad
 B=\frac{D_0b+\ell_b}{4e^2},\quad
 C=\frac{D_0c+\ell_c}{4e^2},\quad
 h=-\frac{D_0^3-fD_0^2+D_0K-L}{8e^3}
\]
とすれば、$u=X/(2e)$等によって同じ等位面を得る。
\emph{これは一次元表現の正規化$A=(1-q)a$とは異なる底変更}である。
同じFricke族の二つの実現を、係数が同一だと解釈してはならない。

\subsection{境界の一方向を除くとPicardの\texorpdfstring{$D_4$}{D4}になる}
\begin{theorem}[整数係数の比較完全列]
\label{thm:lg-integral-picard-sequence}
$U=S\setminus D$が滑らかなら、包含$j:U\hookrightarrow S$は
\begin{equation}\label{eq:lg-integral-sequence}
 0\longrightarrow\mathbb Z\delta_\infty
 \longrightarrow H_2(U,\mathbb Z)
 \xrightarrow{j_*}\langle D_1,D_2,D_3\rangle^\perp
 \longrightarrow0
\end{equation}
を与える。右端は交叉格子として$D_4(-1)$であり、
$H_2(U,\mathbb Z)$はrank $5$の自由加群、$H_1(U,\mathbb Z)=0$である。
交叉形式のradicalは正確に原始部分格子$\mathbb Z\delta_\infty$である。
$\delta_\infty$は向きを除いて境界三角形の$H^1(D,\mathbb Z)$から来る。
\end{theorem}
\begin{proof}
対$(S,D)$のコンパクト台cohomology完全列とPoincar\'e双対性から
\[
\begin{aligned}
 0=H^1(S)&\longrightarrow H^1(D)=\mathbb Z
 \longrightarrow H_c^2(U)\cong H_2(U)\\
 &\longrightarrow H^2(S)\longrightarrow H^2(D)\cong\mathbb Z^3
 \longrightarrow H_c^3(U)\cong H_1(U)\longrightarrow0
\end{aligned}
\]
を得る。三次曲面は有理的なので$H^2(S,\mathbb Z)=\operatorname{Pic}(S)$。
中央の制限写像は三本の$D_i$との交叉数である。
六点blow-upのmarkingを
\[
 D_i=L-E_{2i-1}-E_{2i}\quad(i=1,2,3),\qquad
 L^2=1,\quad E_i^2=-1
\]
と取ると、$E_1,E_3,E_5$の交叉数ベクトルが標準基底になるので全射である。
従って完全列と$H_1(U)=0$が従う。核側・商側が自由なので$H_2(U)$も自由である。
右端の格子の基底は
\[
 \alpha_0=L-E_1-E_3-E_5,\quad
 \alpha_1=E_1-E_2,\quad\alpha_2=E_3-E_4,\quad\alpha_3=E_5-E_6
\]
で、そのGram行列は負の$D_4$ Cartan行列である。
コンパクトな二cycleの交叉数は$S$内に移しても変わらず、
右端の形式は非退化だから、左端がちょうどradicalとなる。
右端がtorsion-freeであることから左端は原始的である。
\end{proof}

\begin{theorem}[局所Milnor格子とPicard格子の実際の同型]
\label{thm:lg-D4-Picard}
$(A,B,C)=(2,2,2)$とし、$h_0=7/2$とする。
$p=(-1,-1,-1)$の十分小さいMilnor球を選び、
$0<|h-h_0|\ll1$に対する局所Milnor fiberを$M_h\subset U_{2,2,2;h}$とする。
このとき二つの包含が誘導する写像は整数交叉格子の同型
\begin{equation}\label{eq:actual-D4-comparison}
 H_2(M_h,\mathbb Z)\xrightarrow{\ \sim\ }
 \langle D_1,D_2,D_3\rangle^\perp\subset\operatorname{Pic}(S_h)
\end{equation}
を与える。従って局所消滅サイクルが式\eqref{eq:lg-integral-sequence}の分裂を与え、
\[
 H_2(U_h,\mathbb Z)\cong H_2(M_h,\mathbb Z)\oplus\mathbb Z\delta_\infty
\]
となる。分裂は指定した退化・Milnor近傍を用いるもので、
全パラメータ上の選択に依存しない標準分裂とは主張しない。
\end{theorem}
\begin{proof}
前節の局所分離式から$p$は$D_4$型曲面特異点である。
他の臨界点$(2,2,2)$の値は$-10$なので、$h_0$の近傍では問題にならない。
ADE特異点のMilnor格子の標準定理により$H_2(M_h)$は$D_4(-1)$、
rank $4$、行列式$4$を持つ。
局所cycleは境界$D$を避けるので、その$S_h$内の像は右端格子に入る。
交叉数が保たれ、局所形式が非退化であるため、この写像は単射である。
両格子のrankは$4$、行列式はともに$4$であり、
包含の指数の平方は行列式の比$1$になる。従って指数は$1$で同型である。
この同型の逆を局所cycleの包含と合成すれば完全列が分裂する。
\end{proof}

この定理は、同じLG等位面の中で\emph{局所Milnorの$D_4$と
三角境界に直交するPicardの$D_4$を幾何学的な包含で結ぶ}。
一般の二つの$D_4$格子に抽象同型がある、というだけの主張ではない。
一方、量子代数の導来圏やFukaya圏の同値はここからは従わない。

\begin{corollary}[外部trialityを検出する判別群]
\label{cor:lg-picard-triality}
対称な族を保つ巡回置換$\rho:(u,v,w)\mapsto(v,w,u)$は、
式\eqref{eq:actual-D4-comparison}を通じ、局所Milnor格子とPicard格子の
両方の判別群$(\mathbb Z/2)^2$の三非零元を巡回する。
従って$\rho$の$\operatorname{Aut}(D_4)/W(D_4)\cong S_3$における像は
位数$3$の外部trialityである。
\end{corollary}
\begin{proof}
$\rho$-不変なMilnor球を選ぶと包含は$\rho$と可換する。
境界格子$B=\langle D_1,D_2,D_3\rangle$のGram行列は
\[
 \begin{pmatrix}-1&1&1\\1&-1&1\\1&1&-1\end{pmatrix}.
\]
これは原始部分格子で、$B^*/B$の三非零元は
$(D_1+D_2)/2,(D_2+D_3)/2,(D_3+D_1)/2$で表される。
境界三本を回す$\rho$はこれらを巡回する。
$\operatorname{Pic}(S)$はunimodularなので、原始直交補のgluingは
$B^*/B$と$B^{\perp*}/B^\perp$の等変同型を与える。
従ってPicardの$D_4$の判別群にも同じ作用がある。
Weyl群は判別群へ自明に作用し、$D_4$では自己同型群のWeyl商が$S_3$だから結論を得る。
\end{proof}

\subsection{affine \texorpdfstring{$D_4$}{D4}、無限遠トーラス、混合Hodge構造}
Goldman--NeumannはFricke曲面の滑らかなfiberの交叉形式を
負のaffine $D_4$として計算している。\footnote{
W. M. Goldman, W. D. Neumann, \emph{Homological action of the modular group on some cubic moduli spaces},
Math. Res. Lett. 12 (2005), 575--591,
\href{https://doi.org/10.4310/MRL.2005.v12.n4.a11}{doi:10.4310/MRL.2005.v12.n4.a11},
Theorems 1--2, Proposition 3, \S3;
\href{https://arxiv.org/abs/math/0402039}{arXiv:math/0402039}.}
四つの外側元$o_1,\ldots,o_4$と中央元$c$に対し
\[
 o_i^2=c^2=-2,\quad o_i\cdot c=1,\quad o_i\cdot o_j=0\ (i\ne j),
 \qquad \delta_\infty=o_1+o_2+o_3+o_4+2c
\]
となる。基底$(c,o_1,o_2,o_3,\delta_\infty)$への変換はunimodularであり、
形式は$D_4(-1)\oplus(0)$となる。
例えばPicard側への格子写像は$c\mapsto\alpha_0$、
$o_i\mapsto\alpha_i$ ($i=1,2,3$)、
$o_4\mapsto-2\alpha_0-\alpha_1-\alpha_2-\alpha_3$で明示できる。
これは基底選択後の表示であり、全パラメータの消滅路を標準化するものではない。

同論文\S3の境界plumbingでは無限遠の三次元多様体は
monodromy $-I$の$T^2$束である。従ってその$H_1$は
$\mathbb Z\oplus(\mathbb Z/2)^2$となる。
判別群の二捩れと、前節のpillowcaseに現れたトーラスの同時反転が再び現れるが、
\emph{三次元のmapping torusと二次元のpillowcaseは同じ空間ではない}。

\begin{proposition}[重み二の四方向と重み四の境界方向]
\label{prop:lg-mixed-hodge}
滑らかな$U$の混合Hodge構造は
\[
 \operatorname{Gr}_2^W H^2(U,\mathbb Q)\cong\mathbb Q(-1)^4,
 \qquad\operatorname{Gr}_3^W H^2(U,\mathbb Q)=0,
 \qquad\operatorname{Gr}_4^W H^2(U,\mathbb Q)\cong\mathbb Q(-2)
\]
を持つ。そのcompact-support Hodge--Deligne多項式は
$E_c(U;x,y)=1+4xy+(xy)^2$、Euler数は$6$である。
ここではgraded pieceを同定し、混合Hodge構造の拡大の分裂は主張しない。
\end{proposition}
\begin{proof}
$D$はsimple normal crossings divisorで、各成分は$\PP^1$、
二重交点は三つである。Deligneの重みspectral sequence\footnote{
P. Deligne, \emph{Th\'eorie de Hodge : II}, Publ. Math. IH\'ES 40 (1971), 5--57,
\href{https://doi.org/10.1007/BF02684692}{doi:10.1007/BF02684692},
\S3.2、Corollary 3.2.13.}
において、重み二は$H^2(S)/\langle[D_1],[D_2],[D_3]\rangle$で次元$4$、
重み三は$\bigoplus_i H^1(D_i)(-1)=0$。
重み四は三角形の辺から頂点へのincidence写像の核で次元$1$、Tate型$(2,2)$である。
また$E_c(S)=1+7xy+(xy)^2$、
$E_c(D)=3(1+xy)-3=3xy$なので、差を取れば後半を得る。
\end{proof}

\subsection{例外群\texorpdfstring{$G_2$}{G2}が二つの異なる衝突をまとめる}
対称部分族$A=B=C=\lambda$ではFricke係数は
$b=-4\lambda,c=8h$である。
Boalch--Palubaの定理\footnote{
P. Boalch, R. Paluba, \emph{Symmetric cubic surfaces and $G_2$ character varieties},
J. Algebraic Geom. 25 (2016), 607--631,
\href{https://doi.org/10.1090/jag/668}{doi:10.1090/jag/668},
\S4, Corollary 11, Remark 12;
\href{https://arxiv.org/abs/1305.6594}{arXiv:1305.6594}.}
を本稿の座標へ移すと次の接続を得る。

\begin{proposition}[対称LG等位面の$G_2$実現（既知定理の適用）]
\label{prop:lg-G2-realization}
$G_2$の共役不変量を$(\alpha,\beta)$とする規約で、係数写像は
\begin{align}\label{eq:lg-G2-parameters}
 \alpha&=8h+2-12\lambda,\notag\\
 \beta&=64\lambda^3+48\lambda^2-96\lambda h-12\lambda-2.
\end{align}
$(\alpha,\beta)$がregular semisimple共役類に対応する開集合では、
$U_{\lambda,\lambda,\lambda;h}$は四穴球面の$G_2$指標多様体の
一つの連結成分として実現される。
ここで最初の三局所monodromyは$G_2$の次元6のsemisimple共役類に属し、
第四が指定されたregular semisimple共役類に属する。
全指標多様体には三つの連結成分があり、\emph{一つのLG等位面と全体を同一視しない}。
\end{proposition}
\begin{proof}
Boalch--Palubaの底写像は
$\alpha=c+2+3b$、$\beta=-b^3+3b^2+3bc+3b-2$である。
$b=-4\lambda,c=8h$を代入すると式\eqref{eq:lg-G2-parameters}を得る。
残りは同論文Corollary 11の適用である。
\end{proof}

同論文のregular locusの記述を本稿の座標で確認するには、
二つのsingular cubic条件
\[
 \lambda^2+2\lambda-2h-1=0,\qquad
 -16\lambda^3-3\lambda^2+12\lambda h+4h^2+2h=0
\]
に加え、余分なconic $16\lambda^2-4\lambda-8h-1=0$を区別する。
滑らかな対称Fricke曲面であっても、このconic上で上記のregular semisimple記述を
無条件に使ってはならない。

\begin{corollary}[同じ$G_2$退化内の$D_4$と$A_2$]
\label{cor:G2-D4-A2}
第四局所monodromyの条件をregular unipotent共役類の閉包へ置き換えると、
Boalch--Palubaの底点は$(\alpha,\beta)=(6,6)$となる。
その二つの成分に対応するLG係数は
\[
 (\lambda,h)=(2,7/2)\quad\text{および}\quad
 (\lambda,h)=(-1/4,1/8)
\]
であり、それぞれ$D_4$型および$A_2$型の特異fiberである。
後者は、一次元表現の$S_5$モノドロミーを検出したsimple foldと同じ係数である。
\end{corollary}
\begin{proof}
同論文Remark 12ではこの底点の二成分は
$(b,c)=(-8,28),(1,1)$で、後者は係数fiberに重複度2で現れる。
実際、$b$を決める三次式は
\[
 b^3+6b^2-15b+8=(b+8)(b-1)^2
\]
である。LGへの変換は上記の二組を与える。
前者の$D_4$は既に示した。後者では$p=(1/2,1/2,1/2)$で
$W=1/8$、Hessianのrankは$2$で、核は対角方向。
対角上に$u=v=w=1/2+t$を代入すると正確に
$W=1/8+t^3$となる。分離補題により局所方程式は
非退化二次形式と非零三次項の和、すなわち$A_2$である。
他の三臨界点は$(5/4,-1,-1)$の順列で、臨界値は$-23/32$だから同fiberにない。
\end{proof}

これは単に$D_4$と$G_2$という名前が一致するという接続ではない。
本稿で別々に現れた\emph{四sheetの$D_4$衝突と二sheetのsimple foldが、
同じ例外群の指標多様体の一つの退化を通じて結ばれる}。
ただし係数fiberの重複度2、局所Milnor数2、Poisson臨界環の二重構造は
異なる概念であり、数字が同じことから同一視しない。
また、Boalch--Palubaの$G_2$実現やGoldman--Neumannのaffine $D_4$は既知結果である。
今回の貢献の位置付けは、本稿のLG・非可換中心・表現被覆の具体式へ移し、
包含による整数格子比較とtriality、異なる衝突の関係を明示したことである。
個別の既出性の網羅調査は未了であり、$G_2$型量子代数との同型、
mirror同値、Fukaya圏の同値、一般H3性を証明したとはしない。

\section{二種類のbraid群、\texorpdfstring{$A_2$}{A2}被覆、有限軌道と境界周期}
\label{sec:g2-braids-periods}
本節では対称LG族
\[
 W_\lambda=uvw-\tfrac12(u^2+v^2+w^2)-\lambda(u+v+w),
 \qquad U_{\lambda,h}=\{W_\lambda=h\}
\]
を用いる。Boalch--Palubaの$G_2$実現で現れる三つの成分を
\emph{底を動かして交換する作用}と、各成分の\emph{内部のHurwitz作用}を区別する。
両方に三本のbraid群$B_3$が現れるが、二つの作用は同じではない。
先行研究の基点はBoalch--Paluba\footnote{
P. Boalch, R. Paluba, \emph{Symmetric cubic surfaces and $G_2$ character varieties},
J. Algebraic Geom. 25 (2016), 607--631,
\href{https://doi.org/10.1090/jag/668}{doi:10.1090/jag/668},
\S4.1, \S5, Corollary 17, \S7, Theorem 20;
\href{https://arxiv.org/abs/1305.6594}{arXiv:1305.6594}.}
の明示的な底写像とbraid同変性である。
以下では既知公式の本稿への適用と、そこから得られる具体的帰結を分ける。

\subsection{三つの成分の被覆は\texorpdfstring{$A_2$}{A2}の普遍三次根被覆}
\begin{theorem}[成分ラベルの$S_3$と二つの異なる分岐条件]
\label{thm:g2-component-A2}
対称Fricke係数$(b,c)=(-4\lambda,8h)$から$G_2$の底への写像
\[
 \alpha=c+2+3b,\qquad
 \beta=-b^3+3b^2+3bc+3b-2
\]
は、次の座標変換によって普遍三次式の根を一つ選ぶ写像となる：
\[
 z=b+2,\quad P=\alpha+3,\quad Q=\beta+6\alpha+12,
 \qquad z^3-3Pz+Q=0.
\]
従ってそのGalois閉包の群は$S_3=W(A_2)$であり、
三つの成分ラベルの幾何学的monodromyも$S_3$である。
底のbranch curveは$4P^3-Q^2=0$で、その補集合の基本群は$B_3$となる。

$\mathfrak d=4(\alpha+3)^3-(\beta+6\alpha+12)^2$と置くと、引戻しは
\begin{equation}\label{eq:g2-branch-pullback}
 \mathfrak d|_{b,c}
 =-(b^2+b-c-1)^2(b^2-8b-4c-16).
\end{equation}
写像自体のramification divisorは$b^2+b-c-1=0$である。
一方、$b^2-8b-4c-16=0$はbranch curveの上に乗る
\emph{一般には非分岐な第三のsheet}を記述する。
\end{theorem}
\begin{proof}
$c$を消去すると
\[
 b^3+6b^2+3(1-\alpha)b+\beta+2=0
\]
となり、上記変換で減次三次式を得る。
$t_1+t_2+t_3=0$の平面を$\mathfrak h$とすると、これは
$\mathfrak h/\langle(23)\rangle\to\mathfrak h/S_3$であり、
$P=-(t_1t_2+t_2t_3+t_3t_1)/3$、$Q=-t_1t_2t_3$、$z=t_1$に対応する。
相異なる三根のunordered configurationから重心を除いた空間が
$\{4P^3-Q^2\ne0\}$なので、基本群は三本braid群である。
$G_2$のregular locusへさらに制限しても、一般関数体のGalois群は変わらない。

三次式の判別式は$27\mathfrak d$であり、底写像のJacobianは
$3(b^2+b-c-1)$となる。代入と因数分解により
式\eqref{eq:g2-branch-pullback}を得る。
これは先行論文\S4.1の判別式の、減次三次式による解釈でもある。
第三sheetが非分岐となるという主張は、二因子の交点を除いた一般点についてのもの。
\end{proof}

\begin{corollary}[三重成分と$3A_1$曲面]
\label{cor:g2-cusp-three-nodes}
成分被覆のcuspは$(\alpha,\beta)=(-3,6)$であり、
その上の唯一の係数は$(b,c)=(-2,1)$、すなわち
$(\lambda,h)=(1/2,1/8)$である。このLG等位面には
$(1/2,-1,-1)$の三順列に、ちょうど三つの$A_1$特異点がある。
\end{corollary}
\begin{proof}
cuspで三次式は$(b+2)^3$となる。
対称LGの臨界方程式の差は、例えば$(v-u)(w+1)=0$。
したがって臨界点は対角上か、$(-1,-1,1-\lambda)$の順列に限る。
$\lambda=1/2$で後三点の臨界値は$1/8$、Hessian行列式は$9/4$。
対角点は$t^2-t=1/2$を満たし、その臨界値は$-3t/2-1/4$なので$1/8$ではない。
射影閉包は無限遠で滑らかだから、これで全特異点を尽くす。
\end{proof}

前節の$(\alpha,\beta)=(6,6)$では三次式は$(b+8)(b-1)^2$であった。
$D_4$成分$(b,c)=(-8,28)$ではramification因子の値は$27\ne0$、
$A_2$成分$(1,1)$ではそれが零となる。
したがって$D_4/A_2$という\emph{曲面特異点の型}と、
非分岐/二重という\emph{成分被覆の分岐}は異なる情報を担う。
cuspでの三重被覆も、曲面が一つの$A_2$になるという意味ではなく、上の$3A_1$を生じる。

\subsection{Hurwitz作用は成分を動かさず、LGとPoisson構造を保つ}
\begin{proposition}[LG座標での二生成元]
\label{prop:lg-hurwitz-explicit}
\begin{align*}
 B_1(u,v,w)&=(u,\,2uv-w-2\lambda,\,v),\\
 B_2(u,v,w)&=(w,\,v,\,2vw-u-2\lambda)
\end{align*}
は多項式自己同型であり、$W_\lambda$と体積形式$du\wedge dv\wedge dw$を保つ。
合成を右から作用させる規約で
\[
 B_1B_2B_1=B_2B_1B_2,\qquad B_1B_2(u,v,w)=(w,u,v),
 \qquad(B_1B_2)^3=1.
\]
従って$B_3$の中心は自明に作用し、作用は$\mathrm{PSL}_2(\mathbb Z)$を経由する。
Jacobian Poisson括弧$\{u,v\}=\partial_w W_\lambda$（巡回）と、
滑らかな等位面のPoincar\'e留数形式も保存される。
\end{proposition}
\begin{proof}
逆写像はそれぞれ
$(u,w,2uw-v-2\lambda)$、$(2uv-w-2\lambda,v,u)$。
代入すると$W_\lambda\circ B_i=W_\lambda$、$\det dB_i=1$とbraid関係を得る。
Jacobian括弧は$W$と体積形式から定義されるので保存され、留数の自然性から後半も従う。
この公式はBoalch--Paluba式(19)に$(x,y,z)=(-2u,-2v,-2w)$、$b=-4\lambda$を入れたもの。
同論文Corollary 17により$G_2$実現ともbraid同変である。
\end{proof}

これらの$B_i$は$b,c$を固定するため、固定$(\alpha,\beta)$の三成分を
\emph{一つずつ保存する}。一方、定理\ref{thm:g2-component-A2}の$B_3$は
底のloopから生じ、成分ラベルに$S_3$を作用させる。
抽象群が同じでも、Hurwitz操作と成分ラベルの解析接続を同一視しない。
また$B_i$の非線型性から、一般に固定した射影三次コンパクト化への正則延長は存在しない。
これは元の非可換H3心のVieta延長障害とも整合するが、圏的対応の構成ではない。

\begin{proposition}[一生成元の反復はChebyshev多項式で線型化される]
\label{prop:lg-twist-chebyshev}
$u\ne1$で$v_*=w_*=\lambda/(u-1)$と置くと、$B_1$は
\[
 \binom{v-v_*}{w-w_*}\longmapsto
 M_u\binom{v-v_*}{w-w_*},\qquad
 M_u=\begin{pmatrix}2u&-1\\1&0\end{pmatrix}
\]
となる。第二種Chebyshev多項式$U_n$を用い、任意の$n\ge1$について
\[
 M_u^n=U_{n-1}(u)M_u-U_{n-2}(u)I,
 \qquad U_{-1}=0
\]
が成り立つ。$u^2\ne1$なら固有値は$\xi,\xi^{-1}$、$u=(\xi+\xi^{-1})/2$であり、
$\xi$が$n$乗根なら$B_1^n$はこの平面上で恒等となる。
\end{proposition}
\begin{proof}
平行移動後の公式を展開し、$M_u^2-2uM_u+I=0$に帰着する。
この漸化式は$U_n=2uU_{n-1}-U_{n-2}$と一致するので帰納法で結論する。
固有値が相異なる範囲では対角化できる。
\end{proof}
これは$u$を固定したconic pencil上の一生成元の線型化であり、
二生成元を同時に線型化したという主張ではない。
実数$|u|<1$では固有値が単位円に乗り、非退化な実conicが存在する場合は
その上の回転として読める。$u=\pm1$の退化と、複数生成元による有限軌道は別に調べる必要がある。

\subsection{七点軌道は\texorpdfstring{$S_7$}{S7}を持つが、位数6048の群ではない}
\begin{proposition}[Klein軌道上の完全対称群]
\label{prop:klein-seven-S7}
$(\lambda,h)=(1/4,0)$、すなわちFricke係数$(b,c)=(-1,0)$を考える。
Fricke座標で順序付き七点集合
\[
 \mathcal O=((0,0,0),(0,0,1),(0,1,0),(1,0,0),
              (1,1,0),(1,0,1),(0,1,1))
\]
は一つのHurwitz軌道であり、二生成元の置換表示は
\[
 p=(1\,3\,7\,2)(4\,5\,6),\qquad
 q=(1\,2\,6\,4)(3\,7\,5).
\]
その像は$S_7$全体である。LG座標では七点の各座標を$-1/2$倍する。
\end{proposition}
\begin{proof}
七点はいずれも$xyz+x^2+y^2+z^2-x-y-z=0$を満たす。
既知のHurwitz公式へ代入して表示した置換を得、作用は推移的である。
右から合成する規約で直接計算すると
\[
 q^{-1}p q^{-2}p q^{-1}p^3=(5\,7).
\]
推移的次数7の群は素数次数なのでprimitiveである。
この互換の共役を辺とするグラフが非連結ならその成分は非自明なblock系となるので、
グラフは連結。連結グラフの辺の互換は全対称群を生成し、結論を得る。
小検証器は二置換から生じる$5040$個も全列挙する。
\end{proof}

この七点軌道自体はBoalch--Paluba\S7の既知軌道である。
同論文Theorem 20は、Fano平面の一点を通る三直線から得られる八元数を用い、
\emph{一つの表現の像}を$G_2(2)'\cong U_3(3)$、位数6048の有限単純群として実現する。
これは\emph{七つの表現同値類をbraidが置換する群}$S_7$（位数5040）とは別物である。
同じ軌道に現れるKleinの位数168の群とも区別する。
ここで得た置換表示を、これらの有限群の新たな同型と解釈してはならない。

この例の$G_2$底は$(\alpha,\beta)=(-1,-1)$で、成分ラベル三次式は
\[
 (b+1)(b^2+5b+1)=0.
\]
従ってKlein成分のほかに$\mathbb Q(\sqrt{21})$上の二つの共役成分がある。
底判別式は$\mathfrak d=7$、他方の$G_2$判別式は$\alpha^2-4\beta-12=-7$なので、
この例はregular semisimple範囲内である。
これは単一の$G_2$fiberの成分構造を記述しており、他二成分にも同じ七点軌道があるとは主張しない。

\subsection{境界周期は固定される}
\begin{theorem}[原始境界cycleの正規化周期]
\label{thm:boundary-residue-period}
一般係数$A,B,C$の滑らかな等位面$U=\{W_{A,B,C}=h\}$上で
\[
 \omega=\operatorname{Res}_{U}
 \frac{du\wedge dv\wedge dw}{W_{A,B,C}-h}
       =\frac{du\wedge dv}{uv-w-C}
\]
と正規化する。境界三角形の節点を囲む小さい実二次元トーラス$T$
の向きを適切に選べば、その類は原始radicalの生成元$\delta_\infty$であり、
\begin{equation}\label{eq:constant-boundary-period}
 \int_{\delta_\infty}\omega=(2\pi i)^2
\end{equation}
となる。対称族のHurwitz作用はこの形式と$\delta_\infty$を固定する。
\end{theorem}
\begin{proof}
射影閉包の$u_{\rm hom}=1$というchartで局所座標を$(V,Z,t)$とすると
\[
 F=VZ-\tfrac12t(1+V^2+Z^2)-t^2(A+BV+CZ)-ht^3.
\]
境界節点で$F_t=-1/2$なので$t=\tau(V,Z)$と解ける。
$V=0$または$Z=0$で$t=0$だから$\tau=VZ\,a(V,Z)$、$a(0,0)=2$である。
アフィン座標$(u,v,w)=(1/t,V/t,Z/t)$へ戻して留数形式を計算すると
\[
 \omega=-\frac{dV\wedge dZ}{\tau(V,Z)F_t(V,Z,\tau(V,Z))}
       =g(V,Z)\frac{dV}{V}\wedge\frac{dZ}{Z},\qquad g(0,0)=1.
\]
二変数Cauchy積分公式から小トーラス上の積分は$(2\pi i)^2$。

この小トーラスの類は、境界三角形のdual graphの一周に対応する。
境界の三つの節点のトーラスに対するMayer--Vietoris関係は
向きを揃えると$[T_{12}]=[T_{23}]=[T_{31}]$であり、
これらが対$(S,D)$の完全列の$H^1(D,\mathbb Z)$からの像を生成する。
従って前節の整数完全列の原始生成元$\delta_\infty$を与える。

対称族の$B_i$は$W$と体積形式を保つので$\omega$を保つ。
また$U$の向きを保つ自己同型は交叉形式と原始radicalを保つため、
$B_{i*}\delta_\infty=\pm\delta_\infty$。
式\eqref{eq:constant-boundary-period}が非零なので符号は$+$である。
\end{proof}

従ってbraid作用は$H_2(U,\mathbb Z)/\mathbb Z\delta_\infty\cong D_4(-1)$へ降りる。
ただしこの事実は四つの内部周期が定数であることを意味せず、
混合Hodge構造の拡大の分裂や全Gauss--Manin系の計算でもない。
境界周期が一つ明示的に正規化されたので、次に調べるべき未知量は
内部の四周期と、境界方向を含む拡大の変化に絞られる。

\begin{remark}[既知性と検証範囲]
成分底写像・Hurwitz公式・Klein七点軌道・八元数による有限群の実現は先行研究にある。
本節はそれらの正規化、二種類のbraid作用の分離、三次根被覆の$A_2$解釈、
七点上の置換群、境界周期を一つのLG模型で整理する。
個別結果の文献上の新規性は未確定である。
恒等式、Chebyshev漸化式の有限次数検算、七点の置換群、境界の二重留数は
\texttt{computations/verify\_g2\_braids\_periods.py}で厳密に再現できる。
全次数の漸化式、configuration spaceの位相、境界トーラスの原始性と積分は
本文の通常の証明であり、計算証明書がそれらを形式検証したとはしない。
\end{remark}

\section{conic matching周期、対数座標、非分裂な混合Hodge構造}
\label{sec:logarithmic-periods}
本節では$\mathbb C$上で
\[
 W=uvw-\tfrac12(u^2+v^2+w^2)-Au-Bv-Cw,\qquad U=\{W=h\},
 \qquad\omega=\frac{du\wedge dv}{uv-w-C}
\]
を用いる。境界cycleの向きは
$\Pi_\infty=\int_{\delta_\infty}\omega=(2\pi i)^2$と固定する。
前節で正規化した境界周期に対し、ここでは実際の内部周期を計算する。
適切なconic matching球面では周期が対数に帰着し、係数零の一変数族では
Picard--Fuchs方程式、無限二面体monodromy、代数的な指数化周期を明示できる。
最後に$h=-1/10$を用い、混合Hodge拡大が有理係数でも非分裂となる例を与える。

\subsection{conic pencilの二重被覆上で留数形式を積に分ける}
\begin{theorem}[matching球面の対数周期]
\label{thm:conic-matching-log-period}
$U$が滑らかであると仮定し、
\[
 d(u)=u^2-1,\qquad
 N(u)=(u^2+2Au+2h)(u^2-1)+B^2+C^2+2uBC
\]
と置く。$N$の相異なる単純零点$u_-,u_+$を、他の零点と$\pm1$を避ける
単純な道$\gamma$で結び、道の近傍で$s(u)=\sqrt{u^2-1}$の枝を選ぶ。
conic fiberの消える円を$\gamma$上でつないだmatching球面$\Sigma_\gamma$に、
以下の公式と整合する向きを入れると
\begin{equation}\label{eq:conic-matching-period}
 \int_{\Sigma_\gamma}\omega
 =2\pi i\int_\gamma\frac{du}{s(u)}
 =2\pi i\bigl[\log(u+s(u))\bigr]_{u_-}^{u_+}.
\end{equation}
対数は道に沿って連続に取る。向き・枝・道を指定せずに周期の一価性を主張しない。
\end{theorem}
\begin{proof}
$u\ne\pm1$で
\[
 v_0=\frac{B+uC}{d},\qquad w_0=\frac{C+uB}{d},\qquad
 V=v-v_0,\quad T=w-w_0
\]
と置くと、直接展開により
\[
 W-h=-\tfrac12(T^2-2uVT+V^2)-\frac{N(u)}{2d(u)}.
\]
$a=T-(u+s)V$、$b=T-(u-s)V$とすれば
\begin{equation}\label{eq:conic-multiplicative-model}
 ab=-\frac{N(u)}{d(u)},\qquad
 \omega=\frac{du}{s(u)}\wedge\frac{da}{a}.
\end{equation}
実際、$b-a=2sV$、$uV-T=-(a+b)/2$なので、fiber方向で
$dV/(uV-T)=da/(sa)$となる。$u$依存の補正項は$du$との外積で消える。

$\gamma$の内部では$a,b\ne0$の双曲線に円を取り、両端でその円を点へ縮める。
単純零点では局所式が$ab=\text{unit}\cdot(u-u_\pm)$なので、
円の円錐は滑らかな二次元円板となり、二端を閉じて球面を得る。
各fiber円で$\int da/a=2\pi i$となる向きにすると、Fubiniの定理から
最初の積分式を得る。$z=u+s$は$u=(z+z^{-1})/2$を満たし、
$du/s=dz/z$なので後半も従う。
\end{proof}

この計算で現れる二重被覆は、有理曲線
$u=(z+z^{-1})/2$であり、一般の楕円曲線の積分ではない。
$u=\pm1$では座標変換が特異となるため、その点を端点とするcycleへ
公式を無条件に流用しない。また本定理だけで全matching球面が整数$D_4$基底をなすとはしない。

\subsection{一つの内部消滅球面を完全に計算する}
\begin{theorem}[Morse球面の正確な周期]
\label{thm:explicit-morse-period}
$A=B=C=0$とし、$-1/2<h<0$、$r=\sqrt{-2h}\in(0,1)$とする。
原点のMorse臨界点へ消える実球面$\delta_0(h)$を、$-r\le u\le r$上で
\begin{align*}
 v&=\frac{\sqrt{r^2-u^2}}{\sqrt{1-u^2}}\cos\theta,\\
 w&=uv+\sqrt{r^2-u^2}\sin\theta,\qquad 0\le\theta\le2\pi
\end{align*}
と表示し、$du\wedge d\theta$に対応する向きを取る。このとき
\begin{equation}\label{eq:arcsine-period}
 \boxed{\Pi_0(h)=\int_{\delta_0(h)}\omega
       =4\pi\arcsin\sqrt{-2h}}.
\end{equation}
$\delta_0(h)$は原始な$A_1$消滅cycleで、自己交叉数は$-2$である。
\end{theorem}
\begin{proof}
代入すると
$v^2-2uvw+w^2=r^2-u^2$であり、$W=-r^2/2=h$となる。
各端点では円が一点へ縮み、近傍では$u$が端点からの距離の二次関数となるので、
これは滑らかな埋込球面である。$r\to0$で原点へ収縮し、
原点のHessianは$-I_3$なので通常の$A_1$消滅球面となる。
標準的なMilnor fiberの交叉形式から自己交叉数は$-2$。
もしこの類が整数倍なら自己交叉数もその整数の平方で割り切れるので、原始的である。

上のパラメータ表示で
\[
 \omega=\frac{du\wedge d\theta}{\sqrt{1-u^2}}
\]
となる。表示上分母と分子が同時に零となる点でも、留数形式の滑らかさにより延長する。
従って$2\pi\int_{-r}^{r}(1-u^2)^{-1/2}du=4\pi\arcsin r$を得る。
この族の臨界点は原点と四つの符号点、臨界値は$0,-1/2$のみなので、
考えた開区間のfiberは滑らかである。
\end{proof}

\begin{proposition}[Picard--Fuchsと無限二面体monodromy]
\label{prop:period-dihedral-monodromy}
$h\in\mathbb C\setminus\{0,-1/2\}$上で解析接続した$\Pi_0$と定数$\Pi_\infty$は
二階方程式
\begin{equation}\label{eq:arcsine-PF}
 2h(1+2h)\frac{d^2y}{dh^2}+(1+4h)\frac{dy}{dh}=0
\end{equation}
の独立な二解である。基点を$(-1/2,0)$に取り、実区間に沿う標準的な接続で
$0,-1/2$の小周回を選ぶと、解の列$(\Pi_\infty,\Pi_0)^{\mathsf t}$への作用は
\[
 T_0=\begin{pmatrix}1&0\\0&-1\end{pmatrix},\qquad
 T_{-1/2}=\begin{pmatrix}1&0\\-1&-1\end{pmatrix}.
\]
その像は無限二面体群であり、定数解の部分局所系を持つこのrank $2$の局所系は非分裂である。
これは全rank $5$ Gauss--Manin系の行列を決定したという主張ではない。
\end{proposition}
\begin{proof}
$h=-r^2/2$を使って式\eqref{eq:arcsine-period}を二回微分すれば
式\eqref{eq:arcsine-PF}を得る。$0$の周りでは$r$の符号が変わり、$\Pi_0\mapsto-\Pi_0$。
$-1/2$の周りでは$\arcsin r\mapsto\pi-\arcsin r$だから
$\Pi_0\mapsto4\pi^2-\Pi_0=-\Pi_\infty-\Pi_0$となる。
両行列はinvolutionで、積は
\[
 T_{-1/2}T_0=\begin{pmatrix}1&0\\-1&1\end{pmatrix}
\]
という非自明unipotentであり、無限二面体群を生成する。
定数解と相補的な共通不変直線が$\Pi_0+c\Pi_\infty$で生成されるとすると、
$T_0$から$c=0$、$T_{-1/2}$から$c=1/2$が必要となり矛盾する。
\end{proof}

このmonodromyは$h$を動かす\emph{係数族の解析接続}であり、
前節の\emph{固定fiber内のHurwitz自己同型}とは異なる。
定数を含む解の二次元空間を計算したのであって、全ホモロジーの二次元不変部分を
無条件に同定したものでもない。

\subsection{指数化周期は代数的トーラス座標になる}
滑らかな三次曲面と境界三角形の組$(S,D)$について、整数比較完全列を使い
\[
 \mathcal P(\bar\gamma)=
 \exp\left(\frac{1}{2\pi i}\int_\gamma\omega\right),\qquad
 \bar\gamma\in H_2(U,\mathbb Z)/\mathbb Z\delta_\infty
                  \cong D^\perp
\]
と置く。境界周期が$(2\pi i)^2$なので持上げによらない。
向きを整合させれば、これは線束の制限
$D^\perp\to\operatorname{Pic}^0(D)\cong\mathbb C^*$と一致する。
この比較はFriedmanの既知のperiod定理\footnote{
R. Friedman, \emph{On the geometry of anticanonical pairs},
\href{https://arxiv.org/abs/1502.02560v2}{arXiv:1502.02560v2} (2016),
\S3, Proposition 3.12, pp. 21--23.
同原稿では境界積分が1となる形式を用いるため、ここでは$\omega/(2\pi i)^2$を代入する。}
による。Gross--Hacking--Keel\footnote{
M. Gross, P. Hacking, S. Keel, \emph{Moduli of surfaces with an anti-canonical cycle},
Compos. Math. 151 (2015), 265--291,
\href{https://doi.org/10.1112/S0010437X14007611}{doi:10.1112/S0010437X14007611},
Definition 1.5, Theorem 1.8.}
のLooijenga pairのperiod mapと同じ対象である。
従って本節の積分は、単なる特殊関数の表示ではなく、
三角境界への線束制限と混合Hodge拡大の具体的な値を与える。

\begin{corollary}[代数的なmatching periodと一次元の明示式]
\label{cor:algebraic-exponential-period}
定理\ref{thm:conic-matching-log-period}のcycleについて
\[
 \mathcal P(\bar\Sigma_\gamma)=\frac{u_++s(u_+)}{u_-+s(u_-)}
\]
となる。従って、零点と平方根を選ぶ有限代数被覆上でこの指数化周期は代数的である。
係数零の族で$\mathcal Q(h)=\mathcal P(\bar\delta_0(h))$と置くと
\begin{equation}\label{eq:algebraic-period-trace}
 \mathcal Q=(\sqrt{1+2h}-i\sqrt{-2h})^2,
 \qquad
 \boxed{\mathcal Q+\mathcal Q^{-1}=2+8h}.
\end{equation}
特に$h=(\mathcal Q-1)^2/(8\mathcal Q)$である。
二つの局所monodromyはどちらも$\mathcal Q\mapsto\mathcal Q^{-1}$を誘導し、
積の対数的translationは指数化で消える。
\end{corollary}
\begin{proof}
一般式は式\eqref{eq:conic-matching-period}を指数化すれば得る。
$r=\sqrt{-2h}$の実区間では
$\mathcal Q=\exp(-2i\arcsin r)=(\sqrt{1-r^2}-ir)^2$。
和を取ると$\mathcal Q+\mathcal Q^{-1}=2-4r^2=2+8h$となる。
境界周期を足しても指数化には影響しないので、monodromyの主張も従う。
\end{proof}

この二重被覆は$\mathbb G_m$の反転商、すなわち$\mathrm{SL}_2$の
固有値とtraceの関係と同型である。
ただし、ここでの$\mathcal Q$は表現の固有値を定義から代入したものではなく、
\emph{消滅球面の周期を指数化して得た量}である。
前節のChebyshev線型化と同じトーラス反転構造が、今度は周期側に現れた。

\subsection{混合Hodge拡大が有理係数でも非分裂となる具体例}
\begin{theorem}[非捩れperiodが検出する非分裂性]
\label{thm:explicit-nonsplit-MHS}
$A=B=C=0$、$h=-1/10$の滑らかなLG等位面では、
\[
 \mathcal P(\bar\delta_0)=\frac{3-4i}{5}
\]
である。この値は1の冪根ではなく、
\[
 0\longrightarrow\mathbb Q(-1)^4
 \longrightarrow H^2(U,\mathbb Q)
 \longrightarrow\mathbb Q(-2)\longrightarrow0
\]
は有理混合Hodge構造として分裂しない。
\end{theorem}
\begin{proof}
$r=1/\sqrt5$を式\eqref{eq:algebraic-period-trace}へ代入すれば表示値を得る。
もしこれが1の冪根なら、その逆数との和も代数的整数である。
しかし和は$6/5$であり、有理数の代数的整数は整数だから矛盾する。

重みのgraded pieceは前節で計算した通りである。
$F^2H^2(U,\mathbb C)$は対数二形式$\omega$で張られる一次元空間。
有理混合Hodge構造が分裂すれば、この直線は$\mathbb Q(-2)$の有理な持上げとなる。
境界周期で割って$\delta_\infty$上の値を1に正規化した
$[\omega]/(2\pi i)^2$は、そのとき有理cohomology classでなければならない。
よって$\int_{\delta_0}\omega/(2\pi i)^2\in\mathbb Q$となり、
その$2\pi i$倍の指数は1の冪根である。上の非捩れ性に矛盾する。
\end{proof}

これはrank $2$の\emph{局所系}の非分裂性とは別の、
\emph{一つのfiberの混合Hodge構造}の非分裂性である。
また一つのperiodの値が1の冪根である場合に、他の三方向も含めた全拡大が分裂するとはしない。

\begin{remark}[mirror・非Archimedean周期への比較]
Looijenga pairのperiodが積分・線束制限・Hodge拡大を結ぶことは既知理論である。
Lai--Zhou\footnote{J. Lai, Y. Zhou, \emph{Mirror symmetry for log Calabi--Yau surfaces II},
Compos. Math. 161 (2025), 2025--2090,
\href{https://doi.org/10.1112/S0010437X25007080}{doi:10.1112/S0010437X25007080},
\S2, Theorem 1.8.}
はmirror族のperiodを用いたモジュライ的同定を行い、
Karwa--Lai\footnote{S. Karwa, J. Lai,
\emph{Non-archimedean periods for log Calabi--Yau surfaces},
\href{https://arxiv.org/abs/2506.01651}{arXiv:2506.01651} (2025),
Introduction equations (1)--(2), Theorem 6.14.}
は適切なgood reductionの下で非Archimedean periodと代数的periodの一致を示す。
本稿の代数的な比$z_+/z_-$はこれらと比較可能な明示例を与えるが、
本稿で非Archimedean SYZ fibrationやmirror同値を構成したという意味ではない。
本節の一般積分式・特定fiberの非分裂性の個別の既出性は未確認である。
全四内部periodの同時基底、全rank $5$ Gauss--Manin接続、量子圏へのperiodの持上げは未解決。
\end{remark}

\section{四つの整数\texorpdfstring{$D_4$}{D4}周期と、円分的な混合Hodge分裂}
\label{sec:four-D4-periods}
本節では$\mathbb C$上で$U=\{W_{A,B,C}=h\}$を考え、
射影閉包を$S$、境界三角形を$D_1+D_2+D_3$とする。
ここで$D_1=\{u_{\rm hom}=t=0\}$、$D_2,D_3$は残る二辺である。
前節の一つのmatching周期から進み、\emph{整数格子の全四方向}の
指数化周期を同時に計算する。ただし以下のconic chartの開条件を課す：
\begin{equation}\label{eq:four-period-open}
 N(u)=(u^2+2Au+2h)(u^2-1)+B^2+C^2+2uBC,
 \qquad \operatorname{Disc}_uN\ne0,\quad B^2-C^2\ne0.
\end{equation}
これより$N$の四根は相異なり$\pm1$と異なる。この開集合では$U$と$S$は滑らかである。
$B^2=C^2$でも$S$が滑らかな場合はあるので、これは全smooth locusの必要条件ではない。

\subsection{conicの成分選択が整数標識を与える}
\begin{lemma}[境界二重切断と$\mathbb F_1$への縮約]
\label{lem:conic-F1-marking}
平面束$u_{\rm hom}=ut$の残余conicは$S\to\mathbb P^1_u$を定め、
特異fiberは$N$の四根上および無限遠の$D_2+D_3$である。
有限な標識被覆上で各有限特異fiberの一成分$E_1,\ldots,E_4$と
無限遠の一成分$E_5=D_2$を選び、これら五本を縮約して$\mathbb F_1$を得られる。
$\mathbb F_1$の負sectionとfiberの全引戻しを$s,F$とすると
\begin{gather}\label{eq:F1-boundary-classes}
 s^2=-1,\quad sF=1,\quad F^2=0,\quad E_i^2=-1,\notag\\
 D_1=2s+2F-\sum_{i=1}^5 E_i,\qquad
 D_2=E_5,\qquad D_3=F-E_5.
\end{gather}
$E_i$同士は互いに直交し、$s,F$は全$E_i$と直交する。
\end{lemma}
\begin{proof}
残余conicを$(v,w,t)$で書くと、その対称行列は
\[
 \begin{pmatrix}-1&u&-B\\u&-1&-C\\-B&-C&-u^2-2Au-2h\end{pmatrix}
\]
で行列式が$N(u)$となる。単純根では二直線となり、無限遠では$D_2+D_3$を得る。
$u\ne\pm1$での局所形$ab=-N(u)/(u^2-1)$から、根が単純なら全空間は滑らか。
$u=\pm1$に全空間の臨界点があれば$B=\mp C$が必要なので、開条件がそれを排除する。
無限遠では既に三次閉包の滑らかさを確認している。

三次曲面の直線配置から、$D_1,D_2$と交わらず$D_3$と交わる直線$s$を取れる。
例えば六点blow-upの表示$D_i=L-E_{2i-1}-E_{2i}$では旧$E_5$がその一つである。
$s$はこのconic bundleのsectionである。
各特異fiberの$s$と交わらない成分を縮約すると、
$s^2=-1$を保った相対極小な$\mathbb P^1$束となり、これは$\mathbb F_1$である。
$D_1$の像は$s$と交わらない二重切断だから$2s+2F$となる。
五つの縮約点での重複度は1なので表示した境界類を得る。
必要な直線の選択は、smooth cubic上の有限エタールな直線スキームと
四つの根の順序を選ぶ有限被覆上で行える。
\end{proof}

$D_1$の座標を$z=w/v$とする。向きを合わせるため全ての$z$を逆数に変えてよい。
$E_i\cap D_1=p_i$の座標を$z_i$、そのconic fiberの値を$u_i$とすると
\begin{equation}\label{eq:zi-conic-trace}
 z_i+z_i^{-1}=2u_i,\qquad z_i=u_i\pm\sqrt{u_i^2-1}.
\end{equation}
符号は単に四つを任意に取るのではなく、上の縮約に用いた成分で決まる。
他方の成分の座標は$z_i^{-1}$である。

\begin{theorem}[全四方向の代数的周期]
\label{thm:four-integer-D4-periods}
補題の標識の下で、$K=\langle D_1,D_2,D_3\rangle^\perp$は
\[
 \left\{bF+\sum_{i=1}^4 c_iE_i:2b+\sum_i c_i=0\right\}
 \cong\left\{(c_1,\ldots,c_4)\in\mathbb Z^4:\sum_i c_i\equiv0\pmod2\right\}
\]
であり、形式は$-\sum_i c_i^2$となる。
次の四根は\emph{整数}基底をなす：
\[
 \alpha_1=E_1-E_2,\quad\alpha_2=E_2-E_3,\quad
 \alpha_3=E_3-E_4,\quad\alpha_4=E_3+E_4-F.
\]
境界の向きと座標を整合させ、線束制限のperiodを$\mathcal P$と書けば
\begin{equation}\label{eq:four-D4-characters}
 \boxed{\bigl(\mathcal P(\alpha_1),\ldots,\mathcal P(\alpha_4)\bigr)
 =\left(\frac{z_1}{z_2},\frac{z_2}{z_3},\frac{z_3}{z_4},z_3z_4\right).}
\end{equation}
従って$\alpha_i$を持ち上げる整数二cycle$\gamma_i\subset U$について
\[
 \int_{\gamma_i}\omega=2\pi i\log\mathcal P(\alpha_i)
 \pmod{(2\pi i)^2\mathbb Z}.
\]
対数の枝を局所的に選べば、$\delta_\infty,\gamma_1,\ldots,\gamma_4$は
$H_2(U,\mathbb Z)$の基底を与え、その全周期をこの式で記述できる。
これは全パラメータ上の標準cycle基底や、代数的de Rham基底での接続行列を
選択なしに与えるという主張ではない。
\end{theorem}
\begin{proof}
$as+bF+\sum c_iE_i$と$D_2,D_3$との直交性から$c_5=a=0$、
$D_1$との直交性から$2b+\sum_{i=1}^4c_i=0$を得る。
四根の係数行列の行列式は2で、偶数和格子の指数も2なので整数基底となる。
そのGram行列は
\[
 \begin{pmatrix}-2&1&0&0\\1&-2&1&1\\0&1&-2&0\\0&1&0&-2\end{pmatrix}.
\]

$E_i$は$D_1$だけを$p_i$で横断し、一般のconic fiber $F_0$も
$D_1$だけと二点$q,q'$で交わる。その二点の$z$座標の積は1である。
したがって$\xi=bF+\sum c_iE_i$の境界制限は、$D_1$上の次数零因子
$b(q+q')+\sum c_ip_i$と、他の二辺上の自明化で計算できる。
$D_1$の両端を$0,\infty$とすると、この自明化のgluing値は、
向きを固定して
\[
 (z(q)z(q'))^b\prod_i z_i^{c_i}=\prod_i z_i^{c_i}
\]
となる。これを四根へ代入すれば式\eqref{eq:four-D4-characters}を得る。

積分との一致はFriedmanのperiod比較定理\footnote{
R. Friedman, \emph{On the geometry of anticanonical pairs},
\href{https://arxiv.org/abs/1502.02560}{arXiv:1502.02560v2},
Proposition 3.12, pp. 21--23。境界周期1の規約へ$\omega/(2\pi i)^2$で合わせる。}
と既存の整数完全列
$0\to\mathbb Z\delta_\infty\to H_2(U,\mathbb Z)\to K\to0$から従う。
同じ比較の証明のように、曲線の差から境界との交点の小円板を除き、
境界に沿うtubeでつなぐと持上げ$\gamma_i$を具体的に取れる。
商の整数基底を任意に持ち上げれば、原始核の生成元と合わせて全体の基底となる。
\end{proof}

\subsection{root hypertorusと四つの局所座標}
\begin{corollary}[周期torus、Weyl型判別式、局所座標]
\label{cor:conic-period-torus}
$(z_1,z_2,z_3,z_4)\mapsto\mathcal P$は
\[
 (\mathbb C^*)^4\longrightarrow\operatorname{Hom}(D_4,\mathbb C^*)
\]
という次数2のisogenyで、核は同時符号変更
$(z_i)\mapsto(-z_i)$である。
四周期は式\eqref{eq:four-period-open}の有限標識被覆上で局所エタール座標となる。
さらに$N$がmonicである規約では
\begin{equation}\label{eq:conic-D4-discriminant}
 \operatorname{Disc}_uN
 =\frac{\prod_{i<j}(z_i-z_j)^2(z_iz_j-1)^2}
 {2^{12}(z_1z_2z_3z_4)^6}.
\end{equation}
すなわちこのchartでの退化条件は$D_4$のroot character
$z_i/z_j,z_iz_j$が1になる条件である。
\end{corollary}
\begin{proof}
指数行列は
$\left(\begin{smallmatrix}1&-1&0&0\\0&1&-1&0\\0&0&1&-1\\0&0&1&1\end{smallmatrix}\right)$
でdeterminantは2。四periodが1なら全$z_i$が等しくその平方は1となる。
係数からquarticの四係数への写像は
\[
 (A,B,C,h)\longmapsto(2A,2h-1,2BC-2A,B^2+C^2-2h)
\]
で、Jacobian determinantは$16(C^2-B^2)$。
根の順序付けは単純根上でエタール、$du_i/d\log z_i=(z_i-z_i^{-1})/2$も非零。
これと指数行列のdeterminantから四periodの独立性を得る。
最後に
\[
 u_i-u_j=\frac{(z_i-z_j)(z_iz_j-1)}{2z_iz_j}
\]
をVandermondeの積へ代入すれば判別式の公式となる。
\end{proof}

標識三次曲面のperiod空間と$D_4$ root hypertorus補集合との関係は
古典的なNaruki--Looijengaの構造であり、Schockの\S2\footnote{
N. Schock, \emph{The $W(E_6)$-invariant birational geometry of the moduli space of marked cubic surfaces},
Math. Nachr. 297 (2024), 2639--2666,
\href{https://doi.org/10.1002/mana.202300459}{doi:10.1002/mana.202300459}, \S2.}
にも記述されている。本節はそれを新しい一般Torelli定理として主張せず、
本稿のconicの四根から四periodを直接再構成する式として与える。
ここでのWeyl作用はroot格子の標識の変更であり、
固定fiberのHurwitz自己同型や先の五重被覆の$S_5$と同じ作用ではない。

\subsection{有理Hodge分裂の必要十分条件と有理数体上の例}
\begin{theorem}[四周期による有理分裂の完全判定]
\label{thm:four-period-rational-split}
式\eqref{eq:four-period-open}の範囲で、次は同値である。
\begin{enumerate}
 \item $0\to\mathbb Q(-1)^4\to H^2(U,\mathbb Q)\to\mathbb Q(-2)\to0$
 が有理混合Hodge構造として分裂する。
 \item 四つの$\mathcal P(\alpha_i)$が全て1の冪根である。
 \item 四つの$z_i$が全て1の冪根である。
\end{enumerate}
この条件は標識・各平方根の選択に依存しない。
\end{theorem}
\begin{proof}
$F^2H^2(U,\mathbb C)$は$[\omega]$が張る一次元空間である。
拡大が分裂することと、境界周期で正規化した
$[\omega]/(2\pi i)^2$が有理cohomology classであることは同値である。
一方向は$\mathbb Q(-2)$の有理な持上げから従う。
逆にその直線が有理なら、重み4の商へ1を送る有理ベクトルが
Hodge filtrationも保つsectionを与える。
整数基底$\delta_\infty,\gamma_1,\ldots,\gamma_4$上で評価すると、
この有理性は全ての$\log\mathcal P(\alpha_i)/(2\pi i)$の有理性、
すなわち(2)と同値である。
(3)から(2)は明らかで、逆は
$z_3^2=\mathcal P(\alpha_3)\mathcal P(\alpha_4)$から
$z_3,z_4,z_2,z_1$の順に従う。根の置換と各逆数化は捩れ性を変えない。
\end{proof}

\begin{proposition}[分裂・非分裂をともに有理係数で実現する]
\label{prop:rational-cyclotomic-examples}
次の二つのLG等位面はともに有理数体上で定義された滑らかな曲面である。
\begin{align*}
 (A,B,C,h)_{\rm split}&=(1/4,1/2,1/4,1/8),\\
 (A,B,C,h)_{\rm nonsplit}&=(-65/16,75/16,-65/16,3575/288).
\end{align*}
前者の$H^2$は有理混合Hodge構造として分裂し、後者は分裂しない。
\end{proposition}
\begin{proof}
前者では
\[
 N(u)=\frac{(2u+1)(8u^3-6u+1)}{16},\qquad
 16z^4N\!\left(\frac{z+z^{-1}}2\right)=\frac{z^9-1}{z-1}.
\]
四根は$\cos(2\pi k/9)$ ($1\le k\le4$)であり、
どの成分選択でも全$z_i$は9乗根となる。
$N$はsquare-freeで、$N(1)=9/16,N(-1)=1/16$だから開条件を満たし、前定理が適用される。
後者では
\[
 N(u)=(u-5/4)(u-5/3)(u-17/8)(u-37/12).
\]
これは$z_i$をそれぞれ$2,3,4,6$またはその逆数から選ぶ場合であり、
全て非捩れである。根は相異なり$\pm1$と異なるので滑らかで、前定理から非分裂となる。
\end{proof}

分裂例が示すのは\emph{有理}混合Hodge構造の分裂である。
この例でも全periodが1というわけではなく、整数係数の拡大の自明性とは異なる。
円分的periodと、元の量子化パラメータ$q^3=1$は別の量であり、
ここから有限Hurwitz軌道や非可換圏の同値は結論しない。

\begin{remark}[現在到達した範囲]
四つの内部方向について、単なる独立なmatching積分ではなく、
整数$D_4$基底・線束制限・指数化積分を同時に記述した。
有理分裂判定はFriedman--Carlsonの既知のperiod解釈の、この座標での具体化である。
一般のroot hypertorus理論や混合Hodgeの分類自体の新規性を主張しない。
本稿での具体式・例の既出性の網羅調査は未了。
全rank五の代数的Gauss--Manin接続行列、Hurwitzとtrialityの全period座標への作用、
$q$変形を含む非可換周期、mirror/Fukaya同値は別の課題として残る。
\end{remark}

\section{trialityの周期isogeny、三つの$G_2$成分、Hodgeに見えない力学}
\label{sec:triality-period-isogeny}
本節は$\mathbb C$上の対称LG族$U_{\lambda,h}=\{W_\lambda=h\}$を扱う。
同じ三次曲面に作用するHurwitz群、$D_4$の標識変更、
$G_2$の三成分を交換する底のmonodromyを区別したうえで、
それらを周期torusの指数3の写像で結ぶ。
乗法的変数$\xi,\eta$は本節の周期座標であり、
元の量子包絡の低次係数$a,b,c$とは異なる。

\subsection{Hurwitzのホモロジー像と固定周期torus}
\begin{proposition}[無限braid作用の有限ホモロジー像]
\label{prop:hurwitz-homology-S3}
滑らかな対称LG族の単連結局所底上で、$H_2(U,\mathbb Z)$を
四つの外側元$o_1,o_2,o_3,o_4$と中央元$c$を持つ負のaffine $D_4$として標識できる。
既存の生成元$B_1,B_2$の作用は、適切に番号を取れば
\[
 (B_1)_*=(o_2\ o_3),\qquad (B_2)_*=(o_1\ o_3),
 \qquad o_4,c\text{ は固定}
\]
となる。従って像は$S_3$、純braid群の像は自明である。
境界radical $\delta=o_1+o_2+o_3+o_4+2c$を割った$K\cong D_4(-1)$では、
$c,o_1,o_2,o_3$が整数基底となる。
これに対する指数化周期を$(\xi,\eta_1,\eta_2,\eta_3)$とすると
\begin{equation}\label{eq:triality-fixed-periods}
 \eta_1=\eta_2=\eta_3=\eta.
\end{equation}
標識を動かす作用は外側三周期の置換であり、固定曲面の周期点はその不動点となる。
\end{proposition}
\begin{proof}
Goldman--Neumannのホモロジー作用の定理\footnote{
W. M. Goldman, W. D. Neumann, \emph{Homological action of the modular group on some cubic moduli spaces},
Math. Res. Lett. 12 (2005), 575--591,
\href{https://doi.org/10.4310/MRL.2005.v12.n4.a11}{doi:10.4310/MRL.2005.v12.n4.a11},
Theorems 1--2, Proposition 3, \S2.}
を本稿の正規化へ移す。$\lambda=0$の四符号臨界点を
$(1,-1,-1),(-1,1,-1),(-1,-1,1),(1,1,1)$と並べると、
$B_1$は第二・第三、$B_2$は第一・第三を交換し、原点は固定する。
体積形式を保つので、同論文の整合した消滅cycleには負号が付かない。
対称smooth底は連結で、射影コンパクト化と境界三角形は局所自明な対の族をなす。
従ってfiberwise自己同型のホモロジー作用を平行移動すれば、同じ行列が得られる。
radicalの$o_4$係数が1なので、残る四元は商の整数基底である。
$B_i^*\omega=\omega$より周期はこれらの置換で不変となり、式を得る。
\end{proof}

\begin{proposition}[不変格子と余不変格子の指数3]
\label{prop:triality-index-three}
$\sigma$を外側三根の巡回置換とし、$s=o_1+o_2+o_3$と置くと
\[
 K^\sigma=\mathbb Zc\oplus\mathbb Zs,
 \qquad K_\sigma=K/(\sigma-1)K=\mathbb Z\bar c\oplus\mathbb Z\bar o.
\]
自然写像は$c\mapsto\bar c,s\mapsto3\bar o$で指数3であり、
\[
 \operatorname{Gram}(c,s)=\begin{pmatrix}-2&3\\3&-6\end{pmatrix}.
\]
従って不動period torusと不変格子の双対torusの間の制限は
\begin{equation}\label{eq:triality-isogeny}
 \mathcal T=\operatorname{Hom}(K_\sigma,\mathbb C^*)\longrightarrow
 \mathcal T_I=\operatorname{Hom}(K^\sigma,\mathbb C^*),
 \qquad(\xi,\eta)\longmapsto(\xi,\eta^3)
\end{equation}
というkernel $\mu_3$のisogenyである。
不変格子と余不変格子を、実ベクトル空間上の折り畳みだけで同一視してはならない。
\end{proposition}
\begin{proof}
巡回で不変な三外側係数は等しく、余不変群では三外側元が等しくなる。
いずれも表示した自由加群なので自然写像の行列は$\operatorname{diag}(1,3)$。
交叉数とtorusの式は直接従う。
\end{proof}

この折り畳みで、$D_4$の12正根のcharacterは
\[
 \eta,\ \xi\eta,\ \xi\eta^2\quad(\text{各重複度3}),\qquad
 \xi,\ \xi\eta^3,\ \xi^2\eta^3\quad(\text{各重複度1})
\]
に移る。従って正根因子の積は
\[
 (1-\eta)^3(1-\xi\eta)^3(1-\xi\eta^2)^3
 (1-\xi)(1-\xi\eta^3)(1-\xi^2\eta^3).
\]
ここでの三重性は根の制限の重複度であり、曲面特異点のMilnor数ではない。

\subsection{対称Fricke係数を実際の二つの周期から再構成する}
\begin{theorem}[対称族の周期による有理式]
\label{thm:symmetric-period-reconstruction}
$\xi,\eta\in\mathbb C^*$について
\begin{align*}
 S&=\eta+\eta^{-1}+\xi\eta+(\xi\eta)^{-1}
                  +\xi\eta^2+(\xi\eta^2)^{-1},\\
 L&=\xi+\xi^{-1}+\xi\eta^3+(\xi\eta^3)^{-1}
                  +\xi^2\eta^3+(\xi^2\eta^3)^{-1}
\end{align*}
とし、
\begin{equation}\label{eq:symmetric-period-reconstruction}
 \lambda=(S+2)/4,\qquad h=(L+3S+4)/8
\end{equation}
と置く。上記六root characterが1でなく、
$\eta,\xi\eta,\xi\eta^2$が$-1$でもないという非空開集合では、
$U_{\lambda,h}$は滑らかで、適合したconic標識における四周期は
\[
 (\mathcal P(\alpha_1),\ldots,\mathcal P(\alpha_4))
       =(\eta,\xi,\eta,\eta)
\]
となる。従って本定理の$\xi,\eta$は単なる形式的torus変数でなく、
実際のLooijenga periodとして実現される。
\end{theorem}
\begin{proof}
前節のconic多項式に$A=B=C=\lambda$を入れると
\begin{equation}\label{eq:symmetric-quartic-factor}
 N(u)=(u+1)\prod_{t\in\{\eta,\xi\eta,\xi\eta^2\}}
                     \left(u+\frac{t+t^{-1}}2\right).
\end{equation}
式\eqref{eq:symmetric-period-reconstruction}を用いた直接展開の恒等式である。
仮定より四根は相異なる。有限fiberの一つは$u=-1$であり、
前節の$B^2\ne C^2$というchartには入らないため、ここを別に扱う。
このfiberの二成分は
\[
 v+w=-\lambda\pm\sqrt{\lambda^2+2\lambda-1-2h}
\]
である。根の単純性により二成分は異なり、交点は境界上にある。
射影閉包は境界に沿って滑らかで、他の単純退化fiberでも全空間は滑らか。
従って$U$は滑らかである。

三つの成分の境界座標を順に
$z_1=-\xi\eta^2,z_2=-\xi\eta,z_3=-\eta$と選び、$z_4=-1$とする。
これらの選択が整数標識に適合することを明示的sectionで確認する。
\[
 k=-\xi\eta^2,\qquad m=(k+k^{-1})/2,\qquad
 \ell=-\tfrac12(\eta+1)(\xi\eta+1)
\]
と置けば、直線$v=ku+\ell,w=m$は$U_{\lambda,h}$の射影閉包上にある。
有限根$u_i=(z_i+z_i^{-1})/2$の選択成分は
$w=z_iv-2\lambda z_i/(z_i-1)$なので、sectionとの交点判定の値は
\begin{gather*}
 \frac{(\xi\eta^2-1)(\xi\eta^3-1)(\xi^2\eta^3-1)}{2\xi\eta^2},\\
 \frac{(\xi\eta-1)(\xi\eta^2-1)(\xi^2\eta^3-1)}{2\xi\eta^2},\qquad
 \frac{(\eta-1)(\xi\eta^2-1)(\xi\eta^3-1)}{2\xi\eta^2}.
\end{gather*}
全て非零である。$u=-1$ではsectionと交わらない方を$E_4$に取る。
両成分の境界座標は同じ$-1$なので、これは$z_4$を変えない。
無限遠でsectionは$D_3$と交わるから、$E_5=D_2$を取れる。
前節と同様にsectionと交わらない五成分を縮約すれば$\mathbb F_1$となる。
各$E_i$と$D_1$の交点は横断的であり、境界gluingの計算は
$u_4=-1$でもそのまま成立する。
従って四periodは$(z_1/z_2,z_2/z_3,z_3/z_4,z_3z_4)=(\eta,\xi,\eta,\eta)$。
積分との比較は前節と同じFriedmanの定理による。
\end{proof}

\subsection{指数3の周期写像がBoalch--Palubaの三成分写像になる}
\begin{theorem}[二つのWeyl商の間の可換図式]
\label{thm:G2-period-isogeny-quotient}
$K$の中央反射と三外側反射の積は$\mathcal T$上で
\[
 r_c(\xi,\eta)=(\xi^{-1},\xi\eta),\qquad
 r_o(\xi,\eta)=(\xi\eta^3,\eta^{-1})
\]
を与え、位数12の$W(G_2)$を生成する。
$\mathcal T/W(G_2)\cong\mathbb A^2_{S,L}$であり、
式\eqref{eq:triality-isogeny}は同じWeyl群に同変である。
$\mathcal T_I$側の基本不変量を$\alpha,\beta$とすれば、商の写像は
\begin{equation}\label{eq:period-isogeny-on-quotients}
 (S,L)\longmapsto
 (\alpha,\beta)=(L,S^3-3SL-9S-6L-12).
\end{equation}
Fricke係数$b_F=-4\lambda=-S-2,c_F=8h=L+3S+4$を用いると
\[
 \alpha=c_F+2+3b_F,\qquad
 \beta=-b_F^3+3b_F^2+3b_Fc_F+3b_F-2.
\]
これはBoalch--Palubaの三成分を選ぶ底写像に正確に一致する。
従って同論文の次数3は、ここでは不変格子から余不変格子への指数3として説明される。
\end{theorem}
\begin{proof}
反射の式は$D_4$の交叉数から計算する。
$y_1=\eta,y_2=\xi\eta,y_3=(\xi\eta^2)^{-1}$と置くと$y_1y_2y_3=1$。
二反射の生成群は三つの$y_i$の置換と同時反転である。
$e_1=\sum y_i,e_2=\sum y_i^{-1}$とすれば、その不変環は
$\mathbb C[e_1,e_2]^{e_1\leftrightarrow e_2}
=\mathbb C[S=e_1+e_2,L=e_1e_2-3]$。

不変格子のtorusには$y'_1=\xi,y'_2=\xi\eta^3,
 y'_3=(\xi^2\eta^3)^{-1}$を用いると、
$\alpha=\sum(y'_i+y_i'^{-1})=L$である。
もう一つの基本不変量は
$\beta=\sum_{i\ne j}y'_i/y'_j=\sum_i(y_i^3+y_i^{-3})$。
Newton恒等式によりこれは$S^3-3SL-9S-6L-12$に等しい。
最後に$b_F,c_F$を代入すれば既知の底写像\footnote{
P. Boalch, R. Paluba, \emph{Symmetric cubic surfaces and $G_2$ character varieties},
J. Algebraic Geom. 25 (2016), 607--631,
\href{https://doi.org/10.1090/jag/668}{doi:10.1090/jag/668},
Lemma 6, equation (12), Corollary 11, Proposition 18 and Remark 19;
\href{https://arxiv.org/abs/1305.6594}{arXiv:1305.6594}.}
と一致する。これは元の$G_2$指標多様体を新たに構成したのではなく、
その既知の底写像と本稿のperiodの間に可換図式を与えるものである。
\end{proof}

この図式では$G_2$の長短根が入れ替わる。
単なる$G_2$の群名の一致ではなく、二つの整数格子とその双対torusを区別する必要がある。
表現論的にも、triality固定の$G_2\subset\mathrm{Spin}(8)$で
$8|_{G_2}=1\oplus7$、$28|_{G_2}=14\oplus7\oplus7$となるので、
$\chi_7=1+S$、$\chi_{14}=2+S+L$を用いると
$b_F=-(\chi_7+1)$、$c_F=\chi_{14}+2\chi_7$として同じ係数を読める。
これはBoalch--Paluba Remark 19のfundamental character記述とも整合する。
ここで構成したのはtorusのisogenyであり、標数零の$G_2$群の非自明な中心isogeny、
Langlands双対群の圏同値、量子代数間のisogenyを構成したとはしない。

\subsection{三成分は有理Hodge構造で一致し、整数構造では異なり得る}
\begin{theorem}[三乗根の差と有理混合Hodge構造]
\label{thm:G2-companion-rational-Hodge}
$G_2$のregularな底点を取り、torus持上げと$\zeta^3=1$を選ぶ。
その三つのFricke成分は、上の商図式では
$(\xi,\eta),(\xi,\zeta\eta),(\xi,\zeta^2\eta)$から得られる。
三持上げが定理\ref{thm:symmetric-period-reconstruction}の開集合に入る範囲では、
成分の$H^2(-,\mathbb Q)$は、標識した重みgraded pieceの同一視の下で
混合Hodge構造として同型である。特に有理Hodge分裂の可否は三成分で同時である。
この同型は標識に依存し、曲面の代数的同型や整数Hodge同型を主張しない。
\end{theorem}
\begin{proof}
regular torus点のWeyl安定化群は自明なので、kernel $\mu_3$の三持上げは
商でも相異なる三点を与える。
四周期は$(\eta,\xi,\eta,\eta)$から
$(\zeta^j\eta,\xi,\zeta^j\eta,\zeta^j\eta)$へ変わる。
境界周期で正規化した積分の差は、外側三方向で$j/3$、中央方向で$0$、
さらに枝による整数を加えたものである。
従って$H_2$の有理基底で外側cycleの持上げに$j\delta_\infty/3$を加減すれば、
$F^2$を張る正規化periodが一致する。
重み二部分は全て型$(1,1)$、重み四商は型$(2,2)$なので、
これは全Hodge filtrationと重みを保つ同型となる。
同値に、periodで分類される拡大類の差が3捩れであり、有理化で消える。
本定理は共通の明示開集合上の主張で、globalな標識選択や
全底上の変動Hodge構造の同型を主張しない。
\end{proof}

\begin{corollary}[Klein曲面の7捩れと二つの21捩れコンパニオン]
\label{cor:Klein-seven-twentyone-Hodge}
$\varepsilon$を原始7乗根として$\xi=\eta=\varepsilon$を取ると
$S=L=-1$、$(b_F,c_F)=(-1,0)$、$(\alpha,\beta)=(-1,-1)$となる。
従って既知のKlein三次曲面
\[
 p^2+r^2+s^2+prs-p-r-s=0
\]
の四periodは全て$\varepsilon$に取れる。
その整数混合Hodge拡大類の位数は正確に$7$である。
同じ$G_2$底の他の二成分は
\[
 b_F=\frac{-5\pm\sqrt{21}}2,\qquad c_F=-3-3b_F
\]
で、整数拡大類の位数は正確に$21$となる。
三成分は有理Hodge構造では分裂して同型だが、
Klein成分と他の二成分は整数混合Hodge構造としては同型でない。
\end{corollary}
\begin{proof}
$S,L$はどちらも非自明な7乗根六個の和なので$-1$。
二つの$G_2$判別因子の値は$7,-7$で非零なのでregularである。
全root characterが1でなく、補助開条件も満たす。
成分三次式は$(b_F+1)(b_F^2+5b_F+1)$となる。
三持上げのperiodは中央$\varepsilon$、外側$\zeta^j\varepsilon$である。
$H_2$の整数比較完全列の核は原始的であるため、
整数Hodge拡大類は四つの正規化対数の$\mathbb C/\mathbb Z$類で表される。
その位数は四つの指数化periodの位数の最小公倍数で、$7,21,21$となる。
この位数は整数graded pieceの基底変更でも変わらない。
有理化すればいずれも零になるから、主張を得る。
\end{proof}

Klein七点braid軌道の存在と$G_2$実現自体はBoalch--Paluba\S7の既知結果である。
ここで比較したのは、その曲面のHodge拡大と同じ底にあるコンパニオンであり、
他の二曲面に同じ有限braid軌道があるとは結論しない。

\subsection{周期が捉えない正のエントロピー}
\begin{proposition}[ホモロジー恒等だが次数が指数増大する自己同型]
\label{prop:period-invisible-entropy}
一般の係数$A,B,C,h$の滑らかなLG等位面上で
\[
 \iota_u=(2vw-u-2A,v,w),\quad
 \iota_v=(u,2uw-v-2B,w),\quad
 \iota_w=(u,v,2uv-w-2C)
\]
とし、$g=\iota_v\iota_w\iota_u\iota_w$と置く（右から合成）。
$g$は$H_2(U,\mathbb Z)$と全periodへ恒等に作用するが、
\[
 \lambda_1(g)=3+2\sqrt2,
 \qquad h_{\rm top}(g)=\log(3+2\sqrt2)>0.
\]
特に有理Hodge分裂は、自己同型の力学が有限であることを意味しない。
\end{proposition}
\begin{proof}
各$\iota$は$W$を保ち体積形式を反転するVieta involutionである。
Goldman--Neumann Theorem 2により、その$H_2$作用は$-I$。
四個の積$g$は$H_2$で恒等となり、留数形式も保存する。
対称族では$g=B_1^2B_2^{-2}$である。
Cayley参照fiberのtorus座標で$B_1,B_2$を表す行列は
\[
 \begin{pmatrix}1&0\\-1&1\end{pmatrix},\qquad
 \begin{pmatrix}1&1\\0&1\end{pmatrix},
 \quad M_g=\begin{pmatrix}1&-2\\-2&5\end{pmatrix}.
\]
$M_g$はlevel 2の合同部分群に属し、固有値は$3\pm2\sqrt2$。
Cantat--Lorayの既知の力学的次数・エントロピー定理\footnote{
S. Cantat, F. Loray, \emph{Dynamics on Character Varieties and Malgrange irreducibility of Painlev\'e VI equation},
Ann. Inst. Fourier 59 (2009), 2927--2978,
\href{https://doi.org/10.5802/aif.2512}{doi:10.5802/aif.2512}.
ここでは著者preprint \href{https://arxiv.org/abs/0711.1579}{arXiv:0711.1579v2},
Theorem B, Theorems 4.7--4.8の規約で照合した。}
を適用して結論する。
\end{proof}

アフィン曲面の$H_2$へ自明に作用することと、
射影コンパクト化の双有理写像としての次数増大は両立する。
$g$は固定したコンパクト化の正則自己同型ではないので、
射影多様体のcohomology上のspectral radiusとの単純な比較はできない。
この例は、周期・格子だけでは見えない作用を局所化した圏や境界変更で
捉える必要性を示す。ただし本稿ではその圏的entropyを計算していない。

\begin{remark}[検証と新規性の位置付け]
有限ホモロジー像、trialityの一般理論、$G_2$底写像、正のentropyは先行結果である。
本節では本稿の周期と整合する標識を構成し、指数3の格子写像から既知の三成分写像を
正確に再現して、三成分の有理Hodge比較とKleinの$7/21$捩れを導いた。
これらの具体的帰結の文献上の既出性の網羅確認は未了である。
整数格子、Laurent恒等式、sectionの非交叉、円分体での例、modular行列は
\texttt{verify\_triality\_period\_isogeny.py}で厳密再現する。
ホモロジー作用、線束period比較、Hodge拡大、entropyは本文の通常の証明と
一次資料への依存であって、計算だけで形式検証したとはしない。
\end{remark}


\section{再現可能性と比較の限界}
\path{computations/verify_all.py}は、中心の恒等式、有限平坦被覆の具体式、
有限体分解型、格子行列および周期の代数恒等式を再現する。
有限体の分解型だけで幾何学的モノドロミーを結論せず、
積分の比較と混合Hodge構造には本文で引用した一般定理を用いる。
全rank五の代数的Gauss--Manin接続行列、全底上の標準的VMHS同型、
skein/DAHAの全パラメータ比較および圏的mirror対応は未構成である。
\UsePublicationReferences{AWAvatars,BoundaryGoldmanNeumann,BousseauSkein,CMR,CantatLorayCayley,CompanionBoundary,CompanionE6,CookeLacabanne,DixonElliptope,EG,EOR,FrickeBoalchPaluba,FrickeFriedman,GoldmanToledo,GoldmanTrace,HavlicekPosta,HikamiCluster,IISDiscriminant,MikoszWeberTriality,SchockD4,StacksPi1,TerwilligerDAHA}
\begin{thebibliography}{99}
\small
\linespread{0.94}\selectfont
\setlength{\itemsep}{0pt}
\ifcsname pubref@DeDekenLowen\endcsname
\bibitem{DeDekenLowen} O. De Deken, W. Lowen,
\emph{Abelian and derived deformations in the presence of $\mathbb Z$-generating geometric helices},
J. Noncommut. Geom. 5 (2011), 477--505,
\href{https://doi.org/10.4171/JNCG/83}{doi:10.4171/JNCG/83}, Theorem 4.5, Proposition 4.6, pp. 495--496.
\fi
\ifcsname pubref@StacksPi1\endcsname
\bibitem{StacksPi1} The Stacks Project Authors,
\emph{The Stacks Project},
\href{https://stacks.math.columbia.edu/tag/0BTX}{Tag 0BTX}, Lemma 58.14.3
(arithmetic--geometric fundamental group exact sequence).
\fi
\ifcsname pubref@StacksPerfect\endcsname
\bibitem{StacksPerfect} The Stacks Project Authors,
\emph{The Stacks Project},
\href{https://stacks.math.columbia.edu/tag/0BDK}{Tag 0BDK}, Lemma 36.31.3.
\fi
\ifcsname pubref@EG\endcsname
\bibitem{EG} P. Etingof, V. Ginzburg,
\emph{Noncommutative del Pezzo surfaces and Calabi--Yau algebras},
J. Eur. Math. Soc. 12 (2010), 1371--1416.
\href{https://arxiv.org/abs/0709.3593v4}{arXiv:0709.3593v4}.
Theorem 3.3.2(i),(ii), Proposition 7.3.1, Appendix 9.1。
\fi
\ifcsname pubref@LBS\endcsname
\bibitem{LBS} L. Le Bruyn, S. P. Smith, M. Van den Bergh,
\emph{Central extensions of three dimensional Artin--Schelter regular algebras},
Math. Z. 222 (1996), 171--212, Theorem 2.6, Corollary 2.7。
\fi
\ifcsname pubref@BG\endcsname
\bibitem{BG} A. Braverman, D. Gaitsgory,
\emph{Poincar\'e--Birkhoff--Witt theorem for quadratic algebras of Koszul type},
J. Algebra 181 (1996), 315--328, Theorem 0.5。
\fi
\ifcsname pubref@IS\endcsname
\bibitem{IS} N. Iyudu, S. Shkarin,
\emph{Three dimensional Sklyanin algebras and Gr\"obner bases},
\href{https://arxiv.org/abs/1601.00564}{arXiv:1601.00564}, Theorem 1.1, Corollary 6.1。
\fi
\ifcsname pubref@NS\endcsname
\bibitem{NS} C. Negron, T. Schedler,
\emph{The Hochschild cohomology ring of a global quotient orbifold},
Adv. Math. 364 (2020), 106978.
\href{https://arxiv.org/abs/1809.08715}{arXiv:1809.08715}, Theorem 1.1, Appendix A。
\fi
\ifcsname pubref@FIOS\endcsname
\bibitem{FIOS} E. Faber, C. Ingalls, S. Okawa, M. Satriano,
\emph{On stacky surfaces and noncommutative surfaces},
\href{https://arxiv.org/abs/2206.13359v2}{arXiv:2206.13359v2}, 8 February 2024,
Theorem 1.5, Corollary 1.7。Brauer捩れ付きの比較として使用する。
\fi
\ifcsname pubref@FPSS\endcsname
\bibitem{FPSS} L. Fu, M. Porta, S. Scherotzke, N. Sibilla,
\emph{Hochschild--Kostant--Rosenberg isomorphism for derived Deligne--Mumford stacks},
\href{https://arxiv.org/abs/2509.00501v3}{arXiv:2509.00501v3}, 20 January 2026,
Proposition 6.17。一般root stackの恒等部分への帰着を比較し、$n=2$ に限る後半の式を $n=3$ へ流用しない。
\fi
\ifcsname pubref@TerwilligerDAHA\endcsname
\bibitem{TerwilligerDAHA} P. Terwilliger,
\emph{The universal Askey--Wilson algebra and DAHA of type $(C_1^\vee,C_1)$},
SIGMA 9 (2013), 047, \href{https://doi.org/10.3842/SIGMA.2013.047}{doi:10.3842/SIGMA.2013.047},
Definitions 2.1, 2.2, 2.6; Lemmas 2.3, 2.5, 2.8.
\fi
\ifcsname pubref@GoldmanTrace\endcsname
\bibitem{GoldmanTrace} W. M. Goldman,
\emph{Trace coordinates on Fricke spaces of some simple hyperbolic surfaces},
in \emph{Handbook of Teichm\"uller Theory}, Vol. II (2009), 611--684,
\href{https://doi.org/10.4171/055-1/16}{doi:10.4171/055-1/16},
\S5.2.1, equation (5.2.1.2).
\fi
\ifcsname pubref@GoldmanToledo\endcsname
\bibitem{GoldmanToledo} W. M. Goldman, D. Toledo,
\emph{Affine cubic surfaces and relative $\mathrm{SL}(2)$-character varieties of compact surfaces},
\href{https://arxiv.org/abs/1006.3838}{arXiv:1006.3838}.
Theorem 1: every smooth-at-infinity cubic with generic tritangent plane is a four-holed-sphere character variety.
\fi
\ifcsname pubref@SchockD4\endcsname
\bibitem{SchockD4} N. Schock,
\emph{The $W(E_6)$-invariant birational geometry of the moduli space of marked cubic surfaces},
Math. Nachr. 297 (2024), 2639--2666,
\href{https://doi.org/10.1002/mana.202300459}{doi:10.1002/mana.202300459}.
Section 2: tritangent triangle and its orthogonal $D_4$ root lattice.
\fi
\ifcsname pubref@IISDiscriminant\endcsname
\bibitem{IISDiscriminant} M. Inaba, K. Iwasaki, M.-H. Saito,
\emph{Dynamics of the sixth Painlev\'e equation},
S\'eminaires \& Congr\`es 14 (2006), 103--167,
\href{https://emis.de/ft/47232}{original text}, \S9.1, Definition 9.1.
Fricke判別式をトレースパラメータへ引き戻した場合の二乗因子について参照する。
本稿の$N_5^2M_5$との同一視は主張しない。
\fi
\ifcsname pubref@HavlicekPosta\endcsname
\bibitem{HavlicekPosta} M. Havl\'{\i}\v{c}ek, S. Po\v{s}ta,
\emph{Center of quantum algebra $U'_q(\mathrm{so}_3)$},
J. Math. Phys. 52 (2011), 043521,
\href{https://doi.org/10.1063/1.3579992}{doi:10.1063/1.3579992}.
三乗根を含むroot-of-unity中心の構造について参照し、規約の完全な同一視は追加検証事項。
\fi
\ifcsname pubref@CantatLorayCayley\endcsname
\bibitem{CantatLorayCayley} S. Cantat, F. Loray,
\emph{Dynamics on Character Varieties and Malgrange irreducibility of Painlev\'e VI equation},
Ann. Inst. Fourier 59 (2009), 2927--2978,
\href{https://doi.org/10.5802/aif.2512}{doi:10.5802/aif.2512}, \S1.5.
\fi
\ifcsname pubref@JangBoundary\endcsname
\bibitem{JangBoundary} S. u. Jang,
\emph{Tropical Dynamics of Markov Surfaces},
\href{https://arxiv.org/abs/2306.11357v2}{arXiv:2306.11357v2},
Lemma 4.8 and Theorem 4.10. 取得した2025年12月版の題名と番号を使用する。
\fi
\ifcsname pubref@MikoszWeberTriality\endcsname
\bibitem{MikoszWeberTriality} M. Mikosz, A. Weber,
\emph{Triality, characteristic classes, $D_4$ and $G_2$ singularities},
J. Homotopy Relat. Struct. 10 (2015), 995--1011,
\href{https://doi.org/10.1007/s40062-014-0094-3}{doi:10.1007/s40062-014-0094-3}, Proposition 3, \S5.
\fi
\ifcsname pubref@DixonElliptope\endcsname
\bibitem{DixonElliptope} K. Dixon,
\emph{The multi-moment maps of the nearly K\"ahler $S^3\times S^3$},
Geom. Dedicata 200 (2019), 351--362,
\href{https://doi.org/10.1007/s10711-018-0375-9}{doi:10.1007/s10711-018-0375-9}, Lemma 3.5.
\fi
\ifcsname pubref@HikamiCluster\endcsname
\bibitem{HikamiCluster} K. Hikami,
\emph{Note on character varieties and cluster algebras},
SIGMA 15 (2019), 003,
\href{https://doi.org/10.3842/SIGMA.2019.003}{doi:10.3842/SIGMA.2019.003}.
\fi
\ifcsname pubref@AWAvatars\endcsname
\bibitem{AWAvatars} N. Cramp\'e et al.,
\emph{The Askey--Wilson algebra and its avatars},
J. Phys. A 54 (2021), 063001,
\href{https://doi.org/10.1088/1751-8121/abd783}{doi:10.1088/1751-8121/abd783}.
\fi
\ifcsname pubref@CookeLacabanne\endcsname
\bibitem{CookeLacabanne} J. Cooke, N. Lacabanne,
\emph{Higher rank Askey--Wilson algebras as skein algebras},
Ann. Inst. Fourier 76 (2026), 1055--1117,
\href{https://doi.org/10.5802/aif.3729}{doi:10.5802/aif.3729}.
\fi
\ifcsname pubref@BousseauSkein\endcsname
\bibitem{BousseauSkein} P. Bousseau,
\emph{Strong positivity for the skein algebras of the $4$-punctured sphere and the $1$-punctured torus},
Comm. Math. Phys. 398 (2023), 1--58,
\href{https://doi.org/10.1007/s00220-022-04512-9}{doi:10.1007/s00220-022-04512-9}.
\fi
\ifcsname pubref@EOR\endcsname
\bibitem{EOR} P. Etingof, A. Oblomkov, E. Rains,
\emph{Generalized double affine Hecke algebras of rank 1 and quantized del Pezzo surfaces},
Adv. Math. 212 (2007), 749--796,
\href{https://doi.org/10.1016/j.aim.2006.11.008}{doi:10.1016/j.aim.2006.11.008}.
\fi
\ifcsname pubref@CMR\endcsname
\bibitem{CMR} L. Chekhov, M. Mazzocco, V. Rubtsov,
\emph{Quantised Painlev\'e monodromy manifolds, Sklyanin and Calabi--Yau algebras},
Adv. Math. 376 (2021), 107442,
\href{https://doi.org/10.1016/j.aim.2020.107442}{doi:10.1016/j.aim.2020.107442},
\href{https://arxiv.org/abs/1905.02772}{arXiv:1905.02772}.
原著v2のDefinition 1.1、式(1.5)、Theorem 1.5、Proposition 4.2を参照。
\fi
\ifcsname pubref@Rains\endcsname
\bibitem{Rains} E. M. Rains,
\emph{Filtered deformations of elliptic algebras},
\href{https://arxiv.org/abs/2107.13540v2}{arXiv:2107.13540v2}, 17 February 2022,
Propositions 2.2, 2.3, Example 2.1, Section 6末尾のConjecture 1とRemark。
\fi
\ifcsname pubref@LVdB\endcsname
\bibitem{LVdB} W. Lowen, M. Van den Bergh,
\emph{Hochschild cohomology of abelian categories and ringed spaces},
Adv. Math. 198 (2005), 172--221,
\href{https://arxiv.org/abs/math/0405227}{arXiv:math/0405227}。
\fi
\ifcsname pubref@LowenObject\endcsname
\bibitem{LowenObject} W. Lowen,
\emph{Hochschild cohomology, the characteristic morphism and derived deformations},
Compos. Math. 144 (2008), 1557--1580,
\href{https://arxiv.org/abs/0707.2602}{arXiv:0707.2602}。
\fi
\ifcsname pubref@Repo\endcsname
\bibitem{Repo} \emph{nc-cubic-moduli research repository},
\href{https://github.com/uedakazushi/nc-cubic-moduli}{GitHub main}。
本稿の依存する個別ファイル・ラベル・確認済み状態は研究ノートの入力表に固定する。
\fi
\ifcsname pubref@CompanionE6\endcsname
\bibitem{CompanionE6}\emph{三乗根量子三次曲面の正常形と形式変形},
日本語研究原稿、2026年11月18日、\PublicationVersion。
\href{https://github.com/uedakazushi/nc-cubic-e6/blob/main/papers/nc-cubic-e6.tex}{\texttt{nc-cubic-e6.tex}}。
\fi
\ifcsname pubref@CompanionFricke\endcsname
\bibitem{CompanionFricke}\emph{量子三次曲面のFricke幾何と周期},
日本語研究原稿、2026年11月18日、\PublicationVersion。
\href{https://github.com/uedakazushi/nc-cubic-e6/blob/main/papers/fricke-geometry.tex}{\texttt{fricke-geometry.tex}}。
\fi
\ifcsname pubref@CompanionBoundary\endcsname
\bibitem{CompanionBoundary}\emph{三次曲面のPell力学と境界特異点},
日本語研究原稿、2026年11月18日、\PublicationVersion。
\href{https://github.com/uedakazushi/nc-cubic-e6/blob/main/papers/cubic-boundary-prym.tex}{\texttt{cubic-boundary-prym.tex}}。
\fi
\ifcsname pubref@AOUCubic\endcsname
\bibitem{AOUCubic}
T. Abdelgadir, S. Okawa, K. Ueda,
\emph{Compact moduli of marked noncommutative cubic surfaces},
\href{https://arxiv.org/abs/2404.00175v1}{arXiv:2404.00175v1} (2024).
\S9, equations (9.1)--(9.10), Lemma 9.2, Theorem 9.4.
\fi
\ifcsname pubref@OUASSurfaces\endcsname
\bibitem{OUASSurfaces}
S. Okawa, K. Ueda,
\emph{Marked Artin--Schelter surfaces of del Pezzo types},
\href{https://arxiv.org/abs/2508.04986v1}{arXiv:2508.04986v1} (2025).
Definition 1.1.
\fi
\ifcsname pubref@BoundaryDHKK\endcsname
\bibitem{BoundaryDHKK} G. Dimitrov, F. Haiden, L. Katzarkov, M. Kontsevich,
\emph{Dynamical systems and categories},
\href{https://arxiv.org/abs/1307.8418}{arXiv:1307.8418},
Definitions 2.1, 2.4, Proposition 2.2, Lemma 2.3, Theorem 2.6.
\fi
\ifcsname pubref@BoundaryGoldmanNeumann\endcsname
\bibitem{BoundaryGoldmanNeumann} W. M. Goldman, W. D. Neumann,
\emph{Homological action of the modular group on some cubic moduli spaces},
Math. Res. Lett. 12 (2005), 575--591,
\href{https://doi.org/10.4310/MRL.2005.v12.n4.a11}{doi:10.4310/MRL.2005.v12.n4.a11},
\href{https://arxiv.org/abs/math/0402039}{arXiv:math/0402039},
Theorems 1--2, Proposition 3.
\fi
\ifcsname pubref@BoundaryDyckerhoff\endcsname
\bibitem{BoundaryDyckerhoff} T. Dyckerhoff,
\emph{Compact generators in categories of matrix factorizations},
Duke Math. J. 159 (2011), 223--274,
\href{https://doi.org/10.1215/00127094-1415869}{doi:10.1215/00127094-1415869},
\href{https://arxiv.org/abs/0904.4713v5}{arXiv:0904.4713v5},
Theorem 5.9, Corollary 6.5.
\fi
\ifcsname pubref@BoundaryHuaKeller\endcsname
\bibitem{BoundaryHuaKeller} Z. Hua, B. Keller,
\emph{Cluster categories and rational curves},
Geom. Topol. 28 (2024), 2569--2634,
\href{https://doi.org/10.2140/gt.2024.28.2569}{doi:10.2140/gt.2024.28.2569},
\href{https://arxiv.org/abs/1810.00749v6}{arXiv:1810.00749v6},
Theorem 5.9, pp.36--37.
\fi
\ifcsname pubref@BoundaryMatherYau\endcsname
\bibitem{BoundaryMatherYau} J. N. Mather, S. S.-T. Yau,
\emph{Classification of isolated hypersurface singularities by their moduli algebras},
Invent. Math. 69 (1982), 243--251,
\href{https://doi.org/10.1007/BF01399504}{doi:10.1007/BF01399504}.
\fi
\ifcsname pubref@BoundaryMumford\endcsname
\bibitem{BoundaryMumford} D. Mumford,
\emph{Prym Varieties I}, in \emph{Contributions to Analysis},
Academic Press (1974), 325--350,
\href{https://doi.org/10.1016/B978-0-12-044850-0.50032-0}{doi:10.1016/B978-0-12-044850-0.50032-0},
\S\S1--3.
\fi
\ifcsname pubref@BoundaryMirandaPersson\endcsname
\bibitem{BoundaryMirandaPersson} R. Miranda, U. Persson,
\emph{On extremal rational elliptic surfaces},
Math. Z. 193 (1986), 537--558,
\href{https://doi.org/10.1007/BF01160474}{doi:10.1007/BF01160474},
Table 5.2.
\fi
\ifcsname pubref@BoundaryHowald\endcsname
\bibitem{BoundaryHowald} J. Howald,
\emph{Multiplier Ideals of Sufficiently General Polynomials},
\href{https://arxiv.org/abs/math/0303203}{arXiv:math/0303203},
Theorems 11--12, Corollary 13.
\fi
\ifcsname pubref@BoundaryCasals\endcsname
\bibitem{BoundaryCasals} R. Casals,
\emph{Lagrangian skeleta and plane curve singularities},
J. Fixed Point Theory Appl. 24 (2022), 34,
\href{https://doi.org/10.1007/s11784-022-00939-8}{doi:10.1007/s11784-022-00939-8},
Theorem 1.1, Corollary 1.2, \S2.1.
\fi
\ifcsname pubref@BoundaryCamara\endcsname
\bibitem{BoundaryCamara} L. M. C\^amara,
\emph{On the classification of quasi-homogeneous curves},
\href{https://arxiv.org/abs/1009.1664v1}{arXiv:1009.1664v1},
Theorem 1(iii), \S4.3.
\fi
\ifcsname pubref@FrickeFriedman\endcsname
\bibitem{FrickeFriedman}
R. Friedman, \emph{On the geometry of anticanonical pairs},
\href{https://arxiv.org/abs/1502.02560v2}{arXiv:1502.02560v2} (2016),
Proposition 3.12.
\fi
\ifcsname pubref@FrickeBoalchPaluba\endcsname
\bibitem{FrickeBoalchPaluba}
P. Boalch, R. Paluba, \emph{Symmetric cubic surfaces and $G_2$ character varieties},
J. Algebraic Geom. 25 (2016), 607--631,
\href{https://doi.org/10.1090/jag/668}{doi:10.1090/jag/668},
Lemma 6, Corollary 11, Proposition 18, Remark 19.
\fi
\ifcsname pubref@PymR13\endcsname
\bibitem{PymR13} B. Pym,
\emph{Quantum deformations of projective three-space},
Adv. Math. 281 (2015), 1216--1241,
\href{https://doi.org/10.1016/j.aim.2015.06.005}{doi:10.1016/j.aim.2015.06.005},
\href{https://arxiv.org/abs/1403.6444}{arXiv:1403.6444}, \S3.4, (9).
\fi
\ifcsname pubref@CompanionAmbient\endcsname
\bibitem{CompanionAmbient}
\emph{非可換三次曲面のambient射と$R(1,3)$成分},
\href{https://github.com/uedakazushi/nc-cubic-ambient/blob/main/papers/nc-cubic-ambient.tex}
{nc-cubic-ambient research manuscript},
\texttt{thm:r13-affine}, \texttt{prop:center}, \texttt{sec:surface-fiber}.
引用範囲は次数付き中心拡大の同型商と中心三次ペンシルであり、
AOUからの支配性の追加仮定は用いない。
\fi
\end{thebibliography}

\end{document}
